\chapter{Rekonstruktion aus Nullproduktprofilen}
\hypertarget{nullproducts.proofs}{}
\hypertarget{np.proof}{}
\noindent\NullproductsMainLink{}\enspace\textperiodcentered\enspace\NullproductsReadingLink{}
\par\medskip
Die kanonische Aussage ist
\FormulaRefAuto{FiniteTwoProductZeroReconstruction}[thm] in Band~37.
Hier werden ihre zusätzlich benötigten Schritte vollständig ausgeführt.
Alle Tabellen verwenden ausschließlich die bereits eingeführten
Schlussregeln und die jeweils angegebenen älteren Aussagen.

Wir kürzen in diesem Kapitel ab:
\[
 \begin{gathered}
 P_S=\PowNE S,\quad P_T=\PowNE T,\\
 H_S=\SemigroupWithZeroAlg{S}{\star}{0_S},\quad
 H_T=\SemigroupWithZeroAlg{T}{\diamond}{0_T},\\
 I=\SemigroupIso{\Phi}{P_S,\star_{\mathcal P}}{P_T,\diamond_{\mathcal P}}.
 \end{gathered}
\]
\(S^2,T^2,P_S^2,P_T^2\) sind stets Produktquadrate.
Die Operationen sind durch ihre Träger festgelegt; auch die spätere
Schreibweise \(xy\) beziehungsweise \(AB\) ist nur eine Abkürzung.
\(A\sim_T B\) bezeichnet die in
\FormulaRefAuto{NullProfileEquivalenceDef}[def] eingeführte Profilrelation.

\Needspace{23\baselineskip}
\section{Das Produktquadrat erkennen}
\hypertarget{np.square}{}
\NullproductsProofTarget{NPProductSquareThree}
\Needspace{14\baselineskip}
\FormulaThmDeltaK[Die drei Produktmengen auf der Quellseite]{%
H_S\dsep S^2=\{0_S,c\}\vdash
 P_S^2=\bigl\{\{0_S\},\{c\},\{0_S,c\}\bigr\}
}{NPProductSquareThree}{%
  \DeltaRow{Mengen}{S\dsep0_S\dsep c}
}
\begin{tabproofwide}
  \proofstepwidestar[1]{H_S}{\rA}
  \proofstepwidestar[2]{S^2=\{0_S,c\}}{\rA}
  \proofstepwidestar[1]{\SemigroupAlg{S}{\star}}{\FormulaRefAuto{SemigroupWithZeroBaseAxiom}[ax]{1}}
  \proofstepwidestar[1]{0_S\in S}{\FormulaRefAuto{SemigroupWithZeroCarrierAxiom}[ax]{1}}
  \proofstepwidestar[1]{S\neq\varnothing}{\FormulaRefAuto{a\in A\vdash A\neq\varnothing}[thm]{4}}
  \proofstepwidestar[1,2]{P_S^2=\PowNE{S^2}}{\FormulaRefAuto{PairProductSquarePowerEquality}[thm]{3,5,2}}
  \nopagebreak
  \proofstepwidestar[]{\PowNE{\{0_S,c\}}=\bigl\{\{0_S\},\{c\},\{0_S,c\}\bigr\}}{\FormulaRefAuto{NonemptyPowersetPair}[thm]}
  \nopagebreak
  \proofstepwidestar[1,2]{P_S^2=\bigl\{\{0_S\},\{c\},\{0_S,c\}\bigr\}}{\rIE{2,7,6}}
\end{tabproofwide}
\Needspace{14\baselineskip}
\FormulaThmDeltaK[Nichtleerheit und Endlichkeit der Zielhalbgruppe]{%
\Finite S\dsep H_S\dsep I\vdash \Finite T\land T\neq\varnothing
}{NPFiniteNonemptyTarget}{%
  \DeltaRow{Mengen}{S\dsep T\dsep0_S}
}
\begin{tabproofwide}
  \proofstepwidestar[1]{\Finite S}{\rA}
  \proofstepwidestar[2]{H_S}{\rA}
  \proofstepwidestar[3]{I}{\rA}
  \proofstepwidestar[1]{\Finite{P_S}}{\FormulaRefAuto{FiniteNonemptyPowerSet}[thm]{1}}
  \proofstepwidestar[3]{\Phi:P_S\bij P_T}{\FormulaRefAuto{SemigroupIsoBijectiveAxiom}[ax]{3}}
  \proofstepwidestar[1,3]{\Finite{P_T}}{\FormulaRefAuto{FiniteBijectionTransport}[thm]{4,5}}
  \proofstepwidestar[1,3]{\Finite T}{\FormulaRefAuto{FiniteNonemptyPowerSetReflectsFiniteness}[thm]{6}}
  \proofstepwidestar[2]{0_S\in S}{\FormulaRefAuto{SemigroupWithZeroCarrierAxiom}[ax]{2}}
  \proofstepwidestar[2]{\{0_S\}\in P_S}{\FormulaRefAuto{SingletonInNonemptyPowerset}[thm]{8}}
  \proofstepwidestar[3]{\Phi:P_S\to P_T}{\FormulaRefAuto{SemigroupIsoFunction}[thm]{3}}
  \proofstepwidestar[2,3]{\Phi(\{0_S\})\in P_T}{\FormulaRefAuto{F\colon A\to B\dsep x\in A\vdash F(x)\in B}[thm]{10,9}}
  \proofstepwidestar[2,3]{\begin{aligned}[t]&\Phi(\{0_S\})\subseteq T\\&{}\land\Phi(\{0_S\})\neq\varnothing\end{aligned}}{\FormulaRefAuto{NonemptyPowersetMembership}[thm]{11}}
  \proofstepwidestar[2,3]{\exists t\in\Phi(\{0_S\})}{\FormulaRefAuto{B\neq\varnothing\vdash\exists x\in B}[thm]{\rAEb{12}}}
  \proofstepwidestar[14]{t\in\Phi(\{0_S\})}{\rA}
  \proofstepwidestar[2,3,14]{t\in T}{\FormulaRefAuto{A\subseteq B\dsep x\in A\vdash x\in B}[thm]{\rAEa{12},14}}
  \proofstepwidestar[2,3,14]{T\neq\varnothing}{\FormulaRefAuto{a\in A\vdash A\neq\varnothing}[thm]{15}}
  \nopagebreak
  \proofstepwidestar[2,3]{T\neq\varnothing}{\rEE{13,14,16}}
  \nopagebreak
  \proofstepwidestar[1,2,3]{\Finite T\land T\neq\varnothing}{\rAI{7,17}}
\end{tabproofwide}
\NullproductsProofTarget{NPTwoProductValuesTransport}
\Needspace{14\baselineskip}
\FormulaThmDeltaK[Die Zielhalbgruppe hat ebenfalls genau zwei Produktwerte]{%
\begin{aligned}[t]
&H_S\dsep S^2=\{0_S,c\}\dsep c\neq0_S
 \dsep\SemigroupAlg{T}{\diamond}\dsep T\neq\varnothing\dsep I\\[-2pt]
&\vdash\exists u\,\exists v\,
 (u\neq v\land T^2=\{u,v\}\land P_T^2=\PowNE{T^2})
\end{aligned}
}{NPTwoProductValuesTransport}{%
  \DeltaRow{Mengen}{S\dsep T\dsep0_S\dsep c\dsep u\dsep v}
}
In dieser Tabelle kürzen \(U=\Phi(\{0_S\})\),
\(V=\Phi(\{c\})\), \(W=\Phi(\{0_S,c\})\),
\(D=T^2\) und \(K=P_T^2\) die wiederkehrenden Mengenterme ab.
\begin{tabproofwide}
  \proofstepwidestar[1]{H_S}{\rA}
  \proofstepwidestar[2]{S^2=\{0_S,c\}}{\rA}
  \proofstepwidestar[3]{c\neq0_S}{\rA}
  \proofstepwidestar[4]{\SemigroupAlg{T}{\diamond}}{\rA}
  \proofstepwidestar[5]{T\neq\varnothing}{\rA}
  \proofstepwidestar[6]{I}{\rA}
  \proofstepwidestar[1]{\SemigroupAlg{S}{\star}}{\FormulaRefAuto{SemigroupWithZeroBaseAxiom}[ax]{1}}
  \proofstepwidestar[1]{0_S\in S}{\FormulaRefAuto{SemigroupWithZeroCarrierAxiom}[ax]{1}}
  \proofstepwidestar[1]{S\neq\varnothing}{\FormulaRefAuto{a\in A\vdash A\neq\varnothing}[thm]{8}}
  \proofstepwidestar[1,2]{P_S^2=\bigl\{\{0_S\},\{c\},\{0_S,c\}\bigr\}}{\FormulaRefAuto{NPProductSquareThree}[thm]{1,2}}
  \proofstepwidestar[2]{c\in S^2}{\rIE{\FormulaRefAuto{a=b\vdash b=a}[thm]{2},\FormulaRefAuto{b\in\{a,b\}}[thm]}}
  \proofstepwidestar[1]{S^2\subseteq S}{\FormulaRefAuto{SemigroupSquareSubset}[thm]{7,9}}
  \proofstepwidestar[1,2]{c\in S}{\FormulaRefAuto{A\subseteq B\dsep x\in A\vdash x\in B}[thm]{12,11}}
  \proofstepwidestar[1]{\{0_S\}\in P_S}{\FormulaRefAuto{SingletonInNonemptyPowerset}[thm]{8}}
  \proofstepwidestar[1,2]{\{c\}\in P_S}{\FormulaRefAuto{SingletonInNonemptyPowerset}[thm]{13}}
  \proofstepwidestar[1,2]{\{0_S,c\}\subseteq S}{\FormulaRefAuto{PairSetSubsetFromMembership}[thm]{8,13}}
  \proofstepwidestar[1,2]{\{0_S,c\}\in P_S}{\FormulaRefAuto{NonemptyPowersetMembership}[thm]{\rAI{16,\FormulaRefAuto{\{a,b\}\neq\varnothing}[thm]}}}
  \proofstepwidestar[6]{\Phi:P_S\to P_T}{\FormulaRefAuto{SemigroupIsoFunction}[thm]{6}}
  \proofstepwidestar[1,2,6]{\Phi[\{\{0_S\},\{c\},\{0_S,c\}\}]=\{U,V,W\}}{\FormulaRefAuto{FunctionImageTriple}[thm]{18,\rAI{14,\rAI{15,17}}}}
  \proofstepwidestar[1]{P_S\neq\varnothing}{\FormulaRefAuto{NonemptyPowerCarrierNonempty}[thm]{9}}
  \proofstepwidestar[5]{P_T\neq\varnothing}{\FormulaRefAuto{NonemptyPowerCarrierNonempty}[thm]{5}}
  \proofstepwidestar[1,5,6]{\Phi[P_S^2]=P_T^2}{\FormulaRefAuto{SemigroupIsoSquareImage}[thm]{6,20,21}}
  \proofstepwidestar[1,2,5,6]{K=\{U,V,W\}}{\rIE{\FormulaRefAuto{a=b\vdash b=a}[thm]{10},22,19}}
  \proofstepwidestar[3]{0_S\neq c}{\FormulaRefAuto{a\neq b\vdash b\neq a}[thm]{3}}
  \proofstepwidestar[3]{\{0_S\}\neq\{c\}}{\rAEa{\FormulaRefAuto{PairNonemptySubsetsDistinct}[thm]{24}}}
  \proofstepwidestar[1,2,6]{\{0_S\}=\{c\}\leftrightarrow U=V}{\FormulaRefAuto{SemigroupIsoEqualityReflection}[thm]{6,\rAI{14,15}}}
  \proofstepwidestar[1,2,3,6]{U\neq V}{\FormulaRefAuto{P\leftrightarrow Q,\neg P\vdash\neg Q}[thm]{26,25}}
  \proofstepwidestar[4,5]{\begin{gathered}D\neq\varnothing\land K\subseteq\PowNE D\\{}\land D\in K\land\forall d\in D\,(\{d\}\in K)\end{gathered}}{\FormulaRefAuto{SemigroupSquarePowerFamily}[thm]{4,5}}
  \proofstepwidestar[1,2,3,4,5,6]{\exists u\,\exists v\,(u\neq v\land D=\{u,v\})}{\FormulaRefAuto{SingletonsAndCarrierForcePair}[thm]{\rAEn{28},\rAEn{28},\rAEn{28},\rAEn{28},23,27}}
  \proofstepwidestar[30]{\exists v\,(u\neq v\land D=\{u,v\})}{\rA}
  \proofstepwidestar[31]{u\neq v\land D=\{u,v\}}{\rA}
  \proofstepwidestar[4,5,31]{P_T^2=\PowNE{T^2}}{\FormulaRefAuto{PairProductSquarePowerEquality}[thm]{4,5,\rAEb{31}}}
  \proofstepwidestar[4,5,31]{\begin{aligned}[t]&\exists u\,\exists v\,(u\neq v\land T^2=\{u,v\}\\&{}\land P_T^2=\PowNE{T^2})\end{aligned}}{\rEI{\rEI{\rAI{\rAEa{31},\rAI{\rAEb{31},32}}}}}
  \nopagebreak
  \proofstepwidestar[4,5,30]{\begin{aligned}[t]&\exists u\,\exists v\,(u\neq v\land T^2=\{u,v\}\\&{}\land P_T^2=\PowNE{T^2})\end{aligned}}{\rEE{30,31,33}}
  \nopagebreak
  \proofstepwidestar[1,2,3,4,5,6]{\begin{aligned}[t]&\exists u\,\exists v\,(u\neq v\land T^2=\{u,v\}\\&{}\land P_T^2=\PowNE{T^2})\end{aligned}}{\rEE{29,30,34}}
\end{tabproofwide}

% Spezielle algebraische Vorstufen zur Existenz der Zielnull.
\Needspace{26\baselineskip}
\section{Die Null aus dem zweielementigen Produktquadrat}
Die Schreibweisen \(S^2,T^2\) bezeichnen Produktquadrate, nicht
kartesische Quadrate. Auf Teilträgern sind die Operationen eingeschränkt.

\Needspace{16\baselineskip}
\FormulaThmDeltaK[Eine Null des Produktquadrats absorbiert den ganzen Träger]{\begin{aligned}[t]
 &\SemigroupAlg{T}{\diamond}\dsep T\ne\varnothing\dsep z\in T^2\dsep\\[-2pt]
 &\forall x\in T^2\,(z\diamond x=z\land x\diamond z=z)
 \vdash\SemigroupWithZeroAlg{T}{\diamond}{z}
\end{aligned}}{NPZeroFromSquareAbsorption}{\DeltaRow{Mengen}{S\dsep T\dsep A\dsep B\dsep P\dsep Q\dsep X\dsep Y\dsep Z\dsep K}\DeltaRow{Elemente}{o\dsep a\dsep b\dsep c\dsep d\dsep x\dsep y\dsep z\dsep u\dsep v}\DeltaRow{Funktionen}{\star\dsep\diamond\dsep F\dsep\Phi\dsep G}}
\begin{tabproofwide}
  \proofstepwidestar[1]{\SemigroupAlg{T}{\diamond}}{\rA}
  \proofstepwidestar[2]{T\ne\varnothing}{\rA}
  \proofstepwidestar[3]{z\in T^2}{\rA}
  \proofstepwidestar[4]{\forall x\in T^2\,(z\diamond x=z\land x\diamond z=z)}{\rA}
  \proofstepwidestar[1,2]{T^2\subseteq T}{\FormulaRefAuto{SemigroupSquareSubset}[thm]{1,2}}
  \proofstepwidestar[1,2,3]{z\in T}{\FormulaRefAuto{A\subseteq B\dsep x\in A\vdash x\in B}[thm]{5,3}}
  \proofstepwidestar[3,4]{z\diamond z=z\land z\diamond z=z}{\rRE{\rUE{4},3}}
  \proofstepwidestar[3,4]{z\diamond z=z}{\rAEa{7}}
  \proofstepwidestar[9]{x\in T}{\rA}
  \proofstepwidestar[1,2,3,9]{z\diamond x\in T^2}{\FormulaRefAuto{SemigroupSquareProductMember}[thm]{1,2,6,9}}
  \proofstepwidestar[1,2,3,9]{x\diamond z\in T^2}{\FormulaRefAuto{SemigroupSquareProductMember}[thm]{1,2,9,6}}
  \proofstepwidestar[1,2,3,4,9]{z\diamond(z\diamond x)=z\land(z\diamond x)\diamond z=z}{\rRE{\rUE{4},10}}
  \proofstepwidestar[1,2,3,4,9]{z\diamond(z\diamond x)=z}{\rAEa{12}}
  \proofstepwidestar[1,2,3,9]{(z\diamond z)\diamond x=z\diamond(z\diamond x)}{\FormulaRefAuto{SemigroupAssociativityAxiom}[ax]{1,6,6,9}}
  \proofstepwidestar[1,2,3,4,9]{z\diamond x=z\diamond(z\diamond x)}{\rIE{8,14}}
  \proofstepwidestar[1,2,3,4,9]{z\diamond x=z}{\rIE{13,15}}
  \proofstepwidestar[1,2,3,4,9]{z\diamond(x\diamond z)=z\land(x\diamond z)\diamond z=z}{\rRE{\rUE{4},11}}
  \proofstepwidestar[1,2,3,4,9]{(x\diamond z)\diamond z=z}{\rAEb{17}}
  \proofstepwidestar[1,2,3,9]{(x\diamond z)\diamond z=x\diamond(z\diamond z)}{\FormulaRefAuto{SemigroupAssociativityAxiom}[ax]{1,9,6,6}}
  \proofstepwidestar[1,2,3,4,9]{(x\diamond z)\diamond z=x\diamond z}{\rIE{8,19}}
  \proofstepwidestar[1,2,3,4,9]{z=x\diamond z}{\rIE{18,20}}
  \proofstepwidestar[1,2,3,4,9]{x\diamond z=z}{\FormulaRefAuto{a=b\vdash b=a}[thm]{21}}
  \proofstepwidestar[1,2,3,4]{\forall x\in T\,(z\diamond x=z)}{\rUI{\rRI{9,16}}}
  \nopagebreak
  \proofstepwidestar[1,2,3,4]{\forall x\in T\,(x\diamond z=z)}{\rUI{\rRI{9,22}}}
  \nopagebreak
  \proofstepwidestar[1,2,3,4]{\SemigroupWithZeroAlg{T}{\diamond}{z}}{\FormulaRefAuto{SemigroupWithZeroDef}[def]{1,6,23,24}}
\end{tabproofwide}

\Needspace{16\baselineskip}
\FormulaThmDeltaK[Absorbierendes Element eines kommutativen idempotenten Zweierquadrats]{\begin{aligned}[t]
 &\SemigroupAlg{T}{\diamond}\dsep T\ne\varnothing\dsep T^2=\{u,v\}\dsep\\[-2pt]
 &\forall x\in T^2\,(x\diamond x=x)\dsep
 \forall x\in T^2\,\forall y\in T^2\,(x\diamond y=y\diamond x)\vdash\\[-2pt]
 &u\diamond v\in T^2\land
 \forall x\in T^2\,((u\diamond v)\diamond x=u\diamond v\land
 x\diamond(u\diamond v)=u\diamond v)
\end{aligned}}{NPTwoElementSemilatticeAbsorber}{\DeltaRow{Mengen}{S\dsep T\dsep A\dsep B\dsep P\dsep Q\dsep X\dsep Y\dsep Z\dsep K}\DeltaRow{Elemente}{o\dsep a\dsep b\dsep c\dsep d\dsep x\dsep y\dsep z\dsep u\dsep v}\DeltaRow{Funktionen}{\star\dsep\diamond\dsep F\dsep\Phi\dsep G}}
\begin{tabproofwide}
  \proofstepwidestar[1]{\SemigroupAlg{T}{\diamond}}{\rA}
  \proofstepwidestar[2]{T\ne\varnothing}{\rA}
  \proofstepwidestar[3]{T^2=\{u,v\}}{\rA}
  \proofstepwidestar[4]{\forall x\in T^2\,(x\diamond x=x)}{\rA}
  \proofstepwidestar[5]{\forall x\in T^2\,\forall y\in T^2\,(x\diamond y=y\diamond x)}{\rA}
  \proofstepwidestar[]{u\in\{u,v\}}{\FormulaRefAuto{a\in\{a,b\}}[thm]}
  \proofstepwidestar[]{v\in\{u,v\}}{\FormulaRefAuto{b\in\{a,b\}}[thm]}
  \proofstepwidestar[3]{\{u,v\}=T^2}{\FormulaRefAuto{a=b\vdash b=a}[thm]{3}}
  \proofstepwidestar[3]{u\in T^2}{\rIE{8,6}}
  \proofstepwidestar[3]{v\in T^2}{\rIE{8,7}}
  \proofstepwidestar[1,2]{T^2\subseteq T}{\FormulaRefAuto{SemigroupSquareSubset}[thm]{1,2}}
  \proofstepwidestar[1,2,3]{u\in T}{\FormulaRefAuto{A\subseteq B\dsep x\in A\vdash x\in B}[thm]{11,9}}
  \proofstepwidestar[1,2,3]{v\in T}{\FormulaRefAuto{A\subseteq B\dsep x\in A\vdash x\in B}[thm]{11,10}}
  \proofstepwidestar[1,2,3]{u\diamond v\in T^2}{\FormulaRefAuto{SemigroupSquareProductMember}[thm]{1,2,12,13}}
  \proofstepwidestar[3,4]{u\diamond u=u}{\rRE{\rUE{4},9}}
  \proofstepwidestar[3,4]{v\diamond v=v}{\rRE{\rUE{4},10}}
  \proofstepwidestar[3,5]{v\diamond u=u\diamond v}{\rRE{\rUE{\rRE{\rUE{5},10}},9}}
  \proofstepwidestar[1,2,3]{(u\diamond v)\diamond u=u\diamond(v\diamond u)}{\FormulaRefAuto{SemigroupAssociativityAxiom}[ax]{1,12,13,12}}
  \proofstepwidestar[1,2,3,5]{(u\diamond v)\diamond u=u\diamond(u\diamond v)}{\rIE{17,18}}
  \proofstepwidestar[1,2,3]{(u\diamond u)\diamond v=u\diamond(u\diamond v)}{\FormulaRefAuto{SemigroupAssociativityAxiom}[ax]{1,12,12,13}}
  \proofstepwidestar[1,2,3,4]{u\diamond v=u\diamond(u\diamond v)}{\rIE{15,20}}
  \proofstepwidestar[1,2,3,4]{u\diamond(u\diamond v)=u\diamond v}{\FormulaRefAuto{a=b\vdash b=a}[thm]{21}}
  \proofstepwidestar[1,2,3,4,5]{(u\diamond v)\diamond u=u\diamond v}{\rIE{22,19}}
  \proofstepwidestar[1,2,3]{(u\diamond v)\diamond v=u\diamond(v\diamond v)}{\FormulaRefAuto{SemigroupAssociativityAxiom}[ax]{1,12,13,13}}
  \proofstepwidestar[1,2,3,4]{(u\diamond v)\diamond v=u\diamond v}{\rIE{16,24}}
  \proofstepwidestar[1,2,3,5]{u\diamond(u\diamond v)=(u\diamond v)\diamond u}{\rRE{\rUE{\rRE{\rUE{5},9}},14}}
  \proofstepwidestar[1,2,3,4,5]{u\diamond(u\diamond v)=u\diamond v}{\rIE{23,26}}
  \proofstepwidestar[1,2,3,5]{v\diamond(u\diamond v)=(u\diamond v)\diamond v}{\rRE{\rUE{\rRE{\rUE{5},10}},14}}
  \proofstepwidestar[1,2,3,4,5]{v\diamond(u\diamond v)=u\diamond v}{\rIE{25,28}}
  \proofstepwidestar[1,2,3,4,5]{(u\diamond v)\diamond u=u\diamond v\land u\diamond(u\diamond v)=u\diamond v}{\rAI{23,27}}
  \proofstepwidestar[1,2,3,4,5]{(u\diamond v)\diamond v=u\diamond v\land v\diamond(u\diamond v)=u\diamond v}{\rAI{25,29}}
  \proofstepwidestar[32]{x\in T^2}{\rA}
  \proofstepwidestar[3,32]{x\in\{u,v\}}{\rIE{3,32}}
  \proofstepwidestar[3,32]{x=u\lor x=v}{\FormulaRefAuto{x\in\{A,B\}\eqvdash(x=A\lor x=B)}[thm]{33}}
  \proofstepwidestar[35]{x=u}{\rA}
  \proofstepwidestar[35]{u=x}{\FormulaRefAuto{a=b\vdash b=a}[thm]{35}}
  \proofstepwidestar[1,2,3,4,5,35]{(u\diamond v)\diamond x=u\diamond v\land x\diamond(u\diamond v)=u\diamond v}{\rIE{36,30}}
  \proofstepwidestar[38]{x=v}{\rA}
  \proofstepwidestar[38]{v=x}{\FormulaRefAuto{a=b\vdash b=a}[thm]{38}}
  \proofstepwidestar[1,2,3,4,5,38]{(u\diamond v)\diamond x=u\diamond v\land x\diamond(u\diamond v)=u\diamond v}{\rIE{39,31}}
  \proofstepwidestar[1,2,3,4,5,32]{(u\diamond v)\diamond x=u\diamond v\land x\diamond(u\diamond v)=u\diamond v}{\rOE{34,35,37,38,40}}
  \nopagebreak
  \proofstepwidestar[1,2,3,4,5]{\forall x\in T^2\,((u\diamond v)\diamond x=u\diamond v\land x\diamond(u\diamond v)=u\diamond v)}{\rUI{\rRI{32,41}}}
  \nopagebreak
  \proofstepwidestar[1,2,3,4,5]{\begin{aligned}[t]&u\diamond v\in T^2\land\\&\forall x\in T^2\,((u\diamond v)\diamond x=u\diamond v\\&\qquad{}\land x\diamond(u\diamond v)=u\diamond v)\end{aligned}}{\rAI{14,42}}
\end{tabproofwide}

\Needspace{16\baselineskip}
\FormulaThmDeltaK[Kommutative Punktprodukte geben kommutative Mengenprodukte]{\begin{aligned}[t]
 &\SemigroupAlg{S}{\star}\dsep X,Y\in\PowNE S\dsep\\[-2pt]
 &\forall x\in X\,\forall y\in Y\,(x\star y=y\star x)\vdash\\[-2pt]
 &X\star_{\mathcal P}Y=Y\star_{\mathcal P}X
\end{aligned}}{NPPowerProductCommutes}{\DeltaRow{Mengen}{S\dsep T\dsep A\dsep B\dsep P\dsep Q\dsep X\dsep Y\dsep Z\dsep K}\DeltaRow{Elemente}{o\dsep a\dsep b\dsep c\dsep d\dsep x\dsep y\dsep z\dsep u\dsep v}\DeltaRow{Funktionen}{\star\dsep\diamond\dsep F\dsep\Phi\dsep G}}
\begin{tabproofwide}
  \proofstepwidestar[1]{\SemigroupAlg{S}{\star}}{\rA}
  \proofstepwidestar[2]{X\in\PowNE S}{\rA}
  \proofstepwidestar[3]{Y\in\PowNE S}{\rA}
  \proofstepwidestar[4]{\forall x\in X\,\forall y\in Y\,(x\star y=y\star x)}{\rA}
  \proofstepwidestar[5]{z\in X\star_{\mathcal P}Y}{\rA}
  \proofstepwidestar[1,2,3,5]{\exists x\in X\,\exists y\in Y\,(z=x\star y)}{\FormulaRefAuto{PowerSemigroupProductWitnesses}[thm]{1,2,3,5}}
  \proofstepwidestar[7]{x\in X\land(y\in Y\land z=x\star y)}{\rA}
  \proofstepwidestar[7]{x\in X}{\rAEa{7}}
  \proofstepwidestar[7]{y\in Y}{\rAEa{\rAEb{7}}}
  \proofstepwidestar[7]{z=x\star y}{\rAEb{\rAEb{7}}}
  \proofstepwidestar[4,7]{x\star y=y\star x}{\rRE{\rUE{\rRE{\rUE{4},8}},9}}
  \proofstepwidestar[1,2,3,7]{y\star x\in Y\star_{\mathcal P}X}{\FormulaRefAuto{PowerSemigroupProductContains}[thm]{1,3,2,9,8}}
  \proofstepwidestar[4,7]{z=y\star x}{\rIE{11,10}}
  \proofstepwidestar[4,7]{y\star x=z}{\FormulaRefAuto{a=b\vdash b=a}[thm]{13}}
  \proofstepwidestar[1,2,3,4,7]{z\in Y\star_{\mathcal P}X}{\rIE{14,12}}
  \proofstepwidestar[1,2,3,4,5]{z\in Y\star_{\mathcal P}X}{\rEEStar{6,7,15}}
  \proofstepwidestar[1,2,3,4]{X\star_{\mathcal P}Y\subseteq Y\star_{\mathcal P}X}{\FormulaRefAuto{A\subseteq B\coloneq\forall x\,(x\in A\rightarrow x\in B)}[def]{\rUI{\rRI{5,16}}}}
  \proofstepwidestar[18]{z\in Y\star_{\mathcal P}X}{\rA}
  \proofstepwidestar[1,2,3,18]{\exists y\in Y\,\exists x\in X\,(z=y\star x)}{\FormulaRefAuto{PowerSemigroupProductWitnesses}[thm]{1,3,2,18}}
  \proofstepwidestar[20]{y\in Y\land(x\in X\land z=y\star x)}{\rA}
  \proofstepwidestar[20]{y\in Y}{\rAEa{20}}
  \proofstepwidestar[20]{x\in X}{\rAEa{\rAEb{20}}}
  \proofstepwidestar[20]{z=y\star x}{\rAEb{\rAEb{20}}}
  \proofstepwidestar[4,20]{x\star y=y\star x}{\rRE{\rUE{\rRE{\rUE{4},22}},21}}
  \proofstepwidestar[4,20]{y\star x=x\star y}{\FormulaRefAuto{a=b\vdash b=a}[thm]{24}}
  \proofstepwidestar[1,2,3,20]{x\star y\in X\star_{\mathcal P}Y}{\FormulaRefAuto{PowerSemigroupProductContains}[thm]{1,2,3,22,21}}
  \proofstepwidestar[4,20]{z=x\star y}{\rIE{25,23}}
  \proofstepwidestar[4,20]{x\star y=z}{\FormulaRefAuto{a=b\vdash b=a}[thm]{27}}
  \proofstepwidestar[1,2,3,4,20]{z\in X\star_{\mathcal P}Y}{\rIE{28,26}}
  \proofstepwidestar[1,2,3,4,18]{z\in X\star_{\mathcal P}Y}{\rEEStar{19,20,29}}
  \nopagebreak
  \proofstepwidestar[1,2,3,4]{Y\star_{\mathcal P}X\subseteq X\star_{\mathcal P}Y}{\FormulaRefAuto{A\subseteq B\coloneq\forall x\,(x\in A\rightarrow x\in B)}[def]{\rUI{\rRI{18,30}}}}
  \nopagebreak
  \proofstepwidestar[1,2,3,4]{X\star_{\mathcal P}Y=Y\star_{\mathcal P}X}{\FormulaRefAuto{A\subseteq B\dsep B\subseteq A\vdash A=B}[thm]{17,31}}
\end{tabproofwide}

\Needspace{16\baselineskip}
\FormulaThmDeltaK[Idempotenz einer Teilmenge mit konservativem Punktprodukt]{\begin{aligned}[t]
 &\SemigroupAlg{S}{\star}\dsep X\in\PowNE S\dsep
 \forall x\in X\,(x\star x=x)\dsep\\[-2pt]
 &\forall x\in X\,\forall y\in X\,(x\star y\in\{x,y\})
 \vdash X\star_{\mathcal P}X=X
\end{aligned}}{NPPowerConservativeIdempotence}{\DeltaRow{Mengen}{S\dsep T\dsep A\dsep B\dsep P\dsep Q\dsep X\dsep Y\dsep Z\dsep K}\DeltaRow{Elemente}{o\dsep a\dsep b\dsep c\dsep d\dsep x\dsep y\dsep z\dsep u\dsep v}\DeltaRow{Funktionen}{\star\dsep\diamond\dsep F\dsep\Phi\dsep G}}
\begin{tabproofwide}
  \proofstepwidestar[1]{\SemigroupAlg{S}{\star}}{\rA}
  \proofstepwidestar[2]{X\in\PowNE S}{\rA}
  \proofstepwidestar[3]{\forall x\in X\,(x\star x=x)}{\rA}
  \proofstepwidestar[4]{\forall x\in X\,\forall y\in X\,(x\star y\in\{x,y\})}{\rA}
  \proofstepwidestar[5]{z\in X\star_{\mathcal P}X}{\rA}
  \proofstepwidestar[1,2,5]{\exists x\in X\,\exists y\in X\,(z=x\star y)}{\FormulaRefAuto{PowerSemigroupProductWitnesses}[thm]{1,2,2,5}}
  \proofstepwidestar[7]{x\in X\land(y\in X\land z=x\star y)}{\rA}
  \proofstepwidestar[7]{x\in X}{\rAEa{7}}
  \proofstepwidestar[7]{y\in X}{\rAEa{\rAEb{7}}}
  \proofstepwidestar[7]{z=x\star y}{\rAEb{\rAEb{7}}}
  \proofstepwidestar[4,7]{x\star y\in\{x,y\}}{\rRE{\rUE{\rRE{\rUE{4},8}},9}}
  \proofstepwidestar[7]{\{x,y\}\subseteq X}{\FormulaRefAuto{PairSetSubsetFromMembership}[thm]{8,9}}
  \proofstepwidestar[4,7]{x\star y\in X}{\FormulaRefAuto{A\subseteq B\dsep x\in A\vdash x\in B}[thm]{12,11}}
  \proofstepwidestar[7]{x\star y=z}{\FormulaRefAuto{a=b\vdash b=a}[thm]{10}}
  \proofstepwidestar[4,7]{z\in X}{\rIE{14,13}}
  \proofstepwidestar[1,2,4,5]{z\in X}{\rEEStar{6,7,15}}
  \proofstepwidestar[1,2,4]{X\star_{\mathcal P}X\subseteq X}{\FormulaRefAuto{A\subseteq B\coloneq\forall x\,(x\in A\rightarrow x\in B)}[def]{\rUI{\rRI{5,16}}}}
  \proofstepwidestar[18]{a\in X}{\rA}
  \proofstepwidestar[3,18]{a\star a=a}{\rRE{\rUE{3},18}}
  \proofstepwidestar[1,2,18]{a\star a\in X\star_{\mathcal P}X}{\FormulaRefAuto{PowerSemigroupProductContains}[thm]{1,2,2,18,18}}
  \proofstepwidestar[1,2,3,18]{a\in X\star_{\mathcal P}X}{\rIE{19,20}}
  \nopagebreak
  \proofstepwidestar[1,2,3]{X\subseteq X\star_{\mathcal P}X}{\FormulaRefAuto{A\subseteq B\coloneq\forall x\,(x\in A\rightarrow x\in B)}[def]{\rUI{\rRI{18,21}}}}
  \nopagebreak
  \proofstepwidestar[1,2,3,4]{X\star_{\mathcal P}X=X}{\FormulaRefAuto{A\subseteq B\dsep B\subseteq A\vdash A=B}[thm]{17,22}}
\end{tabproofwide}

\Needspace{16\baselineskip}
\FormulaThmDeltaK[Das Zweierquadrat im Fall eines verschwindenden zweiten Quadrats]{\begin{aligned}[t]
 &\SemigroupWithZeroAlg{S}{\star}{o}\dsep c\in S\dsep c\star c=o\vdash\\[-2pt]
 &\forall x\in\{o,c\}\,\forall y\in\{o,c\}\,(x\star y=o)
\end{aligned}}{NPZeroPairConstant}{\DeltaRow{Mengen}{S\dsep T\dsep A\dsep B\dsep P\dsep Q\dsep X\dsep Y\dsep Z\dsep K}\DeltaRow{Elemente}{o\dsep a\dsep b\dsep c\dsep d\dsep x\dsep y\dsep z\dsep u\dsep v}\DeltaRow{Funktionen}{\star\dsep\diamond\dsep F\dsep\Phi\dsep G}}
\begin{tabproofwide}
  \proofstepwidestar[1]{\SemigroupWithZeroAlg{S}{\star}{o}}{\rA}
  \proofstepwidestar[2]{c\in S}{\rA}
  \proofstepwidestar[3]{c\star c=o}{\rA}
  \proofstepwidestar[1]{\SemigroupAlg{S}{\star}}{\FormulaRefAuto{SemigroupWithZeroBaseAxiom}[ax]{1}}
  \proofstepwidestar[1]{o\in S}{\FormulaRefAuto{SemigroupWithZeroCarrierAxiom}[ax]{1}}
  \proofstepwidestar[1,2]{\{o,c\}\subseteq S}{\FormulaRefAuto{PairSetSubsetFromMembership}[thm]{5,2}}
  \proofstepwidestar[7]{x\in\{o,c\}}{\rA}
  \proofstepwidestar[8]{y\in\{o,c\}}{\rA}
  \proofstepwidestar[1,2,7]{x\in S}{\FormulaRefAuto{A\subseteq B\dsep x\in A\vdash x\in B}[thm]{6,7}}
  \proofstepwidestar[1,2,8]{y\in S}{\FormulaRefAuto{A\subseteq B\dsep x\in A\vdash x\in B}[thm]{6,8}}
  \proofstepwidestar[7]{x=o\lor x=c}{\FormulaRefAuto{x\in\{A,B\}\eqvdash(x=A\lor x=B)}[thm]{7}}
  \proofstepwidestar[8]{y=o\lor y=c}{\FormulaRefAuto{x\in\{A,B\}\eqvdash(x=A\lor x=B)}[thm]{8}}
  \proofstepwidestar[13]{x=o}{\rA}
  \proofstepwidestar[13]{o=x}{\FormulaRefAuto{a=b\vdash b=a}[thm]{13}}
  \proofstepwidestar[1,2,8]{o\star y=o}{\FormulaRefAuto{SemigroupWithZeroLeftAbsorptionAxiom}[ax]{1,10}}
  \proofstepwidestar[1,2,8,13]{x\star y=o}{\rIE{14,15}}
  \proofstepwidestar[17]{x=c}{\rA}
  \proofstepwidestar[18]{y=o}{\rA}
  \proofstepwidestar[18]{o=y}{\FormulaRefAuto{a=b\vdash b=a}[thm]{18}}
  \proofstepwidestar[1,2,7]{x\star o=o}{\FormulaRefAuto{SemigroupWithZeroRightAbsorptionAxiom}[ax]{1,9}}
  \proofstepwidestar[1,2,7,18]{x\star y=o}{\rIE{19,20}}
  \proofstepwidestar[22]{y=c}{\rA}
  \proofstepwidestar[17]{c=x}{\FormulaRefAuto{a=b\vdash b=a}[thm]{17}}
  \proofstepwidestar[22]{c=y}{\FormulaRefAuto{a=b\vdash b=a}[thm]{22}}
  \proofstepwidestar[3,17]{x\star c=o}{\rIE{23,3}}
  \proofstepwidestar[3,17,22]{x\star y=o}{\rIE{24,25}}
  \proofstepwidestar[1,2,3,7,8,17]{x\star y=o}{\rOE{12,18,21,22,26}}
  \nopagebreak
  \proofstepwidestar[1,2,3,7,8]{x\star y=o}{\rOE{11,13,16,17,27}}
  \nopagebreak
  \proofstepwidestar[1,2,3]{\forall x\in\{o,c\}\,\forall y\in\{o,c\}\,(x\star y=o)}{\rUI{\rRI{7,\rUI{\rRI{8,28}}}}}
\end{tabproofwide}

\Needspace{16\baselineskip}
\FormulaThmDeltaK[Kommutativität auf einem Zweierträger mit Null]{\begin{aligned}[t]
 &\SemigroupWithZeroAlg{S}{\star}{o}\dsep c\in S\vdash\\[-2pt]
 &\forall x\in\{o,c\}\,\forall y\in\{o,c\}\,(x\star y=y\star x)
\end{aligned}}{NPZeroPairCommutative}{\DeltaRow{Mengen}{S\dsep T\dsep A\dsep B\dsep P\dsep Q\dsep X\dsep Y\dsep Z\dsep K}\DeltaRow{Elemente}{o\dsep a\dsep b\dsep c\dsep d\dsep x\dsep y\dsep z\dsep u\dsep v}\DeltaRow{Funktionen}{\star\dsep\diamond\dsep F\dsep\Phi\dsep G}}
\begin{tabproofwide}
  \proofstepwidestar[1]{\SemigroupWithZeroAlg{S}{\star}{o}}{\rA}
  \proofstepwidestar[2]{c\in S}{\rA}
  \proofstepwidestar[1]{\SemigroupAlg{S}{\star}}{\FormulaRefAuto{SemigroupWithZeroBaseAxiom}[ax]{1}}
  \proofstepwidestar[1]{o\in S}{\FormulaRefAuto{SemigroupWithZeroCarrierAxiom}[ax]{1}}
  \proofstepwidestar[1,2]{\{o,c\}\subseteq S}{\FormulaRefAuto{PairSetSubsetFromMembership}[thm]{4,2}}
  \proofstepwidestar[6]{x\in\{o,c\}}{\rA}
  \proofstepwidestar[7]{y\in\{o,c\}}{\rA}
  \proofstepwidestar[1,2,6]{x\in S}{\FormulaRefAuto{A\subseteq B\dsep x\in A\vdash x\in B}[thm]{5,6}}
  \proofstepwidestar[1,2,7]{y\in S}{\FormulaRefAuto{A\subseteq B\dsep x\in A\vdash x\in B}[thm]{5,7}}
  \proofstepwidestar[6]{x=o\lor x=c}{\FormulaRefAuto{x\in\{A,B\}\eqvdash(x=A\lor x=B)}[thm]{6}}
  \proofstepwidestar[7]{y=o\lor y=c}{\FormulaRefAuto{x\in\{A,B\}\eqvdash(x=A\lor x=B)}[thm]{7}}
  \proofstepwidestar[12]{x=o}{\rA}
  \proofstepwidestar[12]{o=x}{\FormulaRefAuto{a=b\vdash b=a}[thm]{12}}
  \proofstepwidestar[1,2,7]{o\star y=o}{\FormulaRefAuto{SemigroupWithZeroLeftAbsorptionAxiom}[ax]{1,9}}
  \proofstepwidestar[1,2,7]{y\star o=o}{\FormulaRefAuto{SemigroupWithZeroRightAbsorptionAxiom}[ax]{1,9}}
  \proofstepwidestar[1,2,7]{o=y\star o}{\FormulaRefAuto{a=b\vdash b=a}[thm]{15}}
  \proofstepwidestar[1,2,7]{o\star y=y\star o}{\rIE{16,14}}
  \proofstepwidestar[1,2,7,12]{x\star y=y\star x}{\rIE{13,17}}
  \proofstepwidestar[19]{x=c}{\rA}
  \proofstepwidestar[20]{y=o}{\rA}
  \proofstepwidestar[20]{o=y}{\FormulaRefAuto{a=b\vdash b=a}[thm]{20}}
  \proofstepwidestar[1,2,6]{x\star o=o}{\FormulaRefAuto{SemigroupWithZeroRightAbsorptionAxiom}[ax]{1,8}}
  \proofstepwidestar[1,2,6]{o\star x=o}{\FormulaRefAuto{SemigroupWithZeroLeftAbsorptionAxiom}[ax]{1,8}}
  \proofstepwidestar[1,2,6]{o=o\star x}{\FormulaRefAuto{a=b\vdash b=a}[thm]{23}}
  \proofstepwidestar[1,2,6]{x\star o=o\star x}{\rIE{24,22}}
  \proofstepwidestar[1,2,6,20]{x\star y=y\star x}{\rIE{21,25}}
  \proofstepwidestar[27]{y=c}{\rA}
  \proofstepwidestar[19]{c=x}{\FormulaRefAuto{a=b\vdash b=a}[thm]{19}}
  \proofstepwidestar[19,27]{y=x}{\rIE{28,27}}
  \proofstepwidestar[19,27]{x=y}{\FormulaRefAuto{a=b\vdash b=a}[thm]{29}}
  \proofstepwidestar[]{x\star x=x\star x}{\rII}
  \proofstepwidestar[19,27]{x\star y=y\star x}{\rIE{30,31}}
  \proofstepwidestar[1,2,6,7,19]{x\star y=y\star x}{\rOE{11,20,26,27,32}}
  \nopagebreak
  \proofstepwidestar[1,2,6,7]{x\star y=y\star x}{\rOE{10,12,18,19,33}}
  \nopagebreak
  \proofstepwidestar[1,2]{\forall x\in\{o,c\}\,\forall y\in\{o,c\}\,(x\star y=y\star x)}{\rUI{\rRI{6,\rUI{\rRI{7,34}}}}}
\end{tabproofwide}

\Needspace{16\baselineskip}
\FormulaThmDeltaK[Idempotenz auf dem Zweierträger im nichtverschwindenden Fall]{\SemigroupWithZeroAlg{S}{\star}{o}\dsep c\in S\dsep c\star c=c
 \vdash\forall x\in\{o,c\}\,(x\star x=x)}{NPZeroPairIdempotent}{\DeltaRow{Mengen}{S\dsep T\dsep A\dsep B\dsep P\dsep Q\dsep X\dsep Y\dsep Z\dsep K}\DeltaRow{Elemente}{o\dsep a\dsep b\dsep c\dsep d\dsep x\dsep y\dsep z\dsep u\dsep v}\DeltaRow{Funktionen}{\star\dsep\diamond\dsep F\dsep\Phi\dsep G}}
\begin{tabproofwide}
  \proofstepwidestar[1]{\SemigroupWithZeroAlg{S}{\star}{o}}{\rA}
  \proofstepwidestar[2]{c\in S}{\rA}
  \proofstepwidestar[3]{c\star c=c}{\rA}
  \proofstepwidestar[1]{\SemigroupAlg{S}{\star}}{\FormulaRefAuto{SemigroupWithZeroBaseAxiom}[ax]{1}}
  \proofstepwidestar[1]{o\in S}{\FormulaRefAuto{SemigroupWithZeroCarrierAxiom}[ax]{1}}
  \proofstepwidestar[1,2]{\{o,c\}\subseteq S}{\FormulaRefAuto{PairSetSubsetFromMembership}[thm]{5,2}}
  \proofstepwidestar[1]{o\star o=o}{\FormulaRefAuto{SemigroupWithZeroLeftAbsorptionAxiom}[ax]{1,5}}
  \proofstepwidestar[8]{x\in\{o,c\}}{\rA}
  \proofstepwidestar[8]{x=o\lor x=c}{\FormulaRefAuto{x\in\{A,B\}\eqvdash(x=A\lor x=B)}[thm]{8}}
  \proofstepwidestar[10]{x=o}{\rA}
  \proofstepwidestar[10]{o=x}{\FormulaRefAuto{a=b\vdash b=a}[thm]{10}}
  \proofstepwidestar[1,10]{x\star x=x}{\rIE{11,7}}
  \proofstepwidestar[13]{x=c}{\rA}
  \proofstepwidestar[13]{c=x}{\FormulaRefAuto{a=b\vdash b=a}[thm]{13}}
  \proofstepwidestar[3,13]{x\star x=x}{\rIE{14,3}}
  \nopagebreak
  \proofstepwidestar[1,3,8]{x\star x=x}{\rOE{9,10,12,13,15}}
  \nopagebreak
  \proofstepwidestar[1,3]{\forall x\in\{o,c\}\,(x\star x=x)}{\rUI{\rRI{8,16}}}
\end{tabproofwide}

\Needspace{16\baselineskip}
\FormulaThmDeltaK[Jedes Produkt des idempotenten Zweierträgers ist ein Faktor]{\begin{aligned}[t]
 &\SemigroupWithZeroAlg{S}{\star}{o}\dsep c\in S\dsep c\star c=c\vdash\\[-2pt]
 &\forall x\in\{o,c\}\,\forall y\in\{o,c\}\,(x\star y\in\{x,y\})
\end{aligned}}{NPZeroPairConservative}{\DeltaRow{Mengen}{S\dsep T\dsep A\dsep B\dsep P\dsep Q\dsep X\dsep Y\dsep Z\dsep K}\DeltaRow{Elemente}{o\dsep a\dsep b\dsep c\dsep d\dsep x\dsep y\dsep z\dsep u\dsep v}\DeltaRow{Funktionen}{\star\dsep\diamond\dsep F\dsep\Phi\dsep G}}
\begin{tabproofwide}
  \proofstepwidestar[1]{\SemigroupWithZeroAlg{S}{\star}{o}}{\rA}
  \proofstepwidestar[2]{c\in S}{\rA}
  \proofstepwidestar[3]{c\star c=c}{\rA}
  \proofstepwidestar[1]{\SemigroupAlg{S}{\star}}{\FormulaRefAuto{SemigroupWithZeroBaseAxiom}[ax]{1}}
  \proofstepwidestar[1]{o\in S}{\FormulaRefAuto{SemigroupWithZeroCarrierAxiom}[ax]{1}}
  \proofstepwidestar[1,2]{\{o,c\}\subseteq S}{\FormulaRefAuto{PairSetSubsetFromMembership}[thm]{5,2}}
  \proofstepwidestar[7]{x\in\{o,c\}}{\rA}
  \proofstepwidestar[8]{y\in\{o,c\}}{\rA}
  \proofstepwidestar[1,2,7]{x\in S}{\FormulaRefAuto{A\subseteq B\dsep x\in A\vdash x\in B}[thm]{6,7}}
  \proofstepwidestar[1,2,8]{y\in S}{\FormulaRefAuto{A\subseteq B\dsep x\in A\vdash x\in B}[thm]{6,8}}
  \proofstepwidestar[7]{x=o\lor x=c}{\FormulaRefAuto{x\in\{A,B\}\eqvdash(x=A\lor x=B)}[thm]{7}}
  \proofstepwidestar[8]{y=o\lor y=c}{\FormulaRefAuto{x\in\{A,B\}\eqvdash(x=A\lor x=B)}[thm]{8}}
  \proofstepwidestar[]{x\in\{x,y\}}{\FormulaRefAuto{a\in\{a,b\}}[thm]}
  \proofstepwidestar[]{y\in\{x,y\}}{\FormulaRefAuto{b\in\{a,b\}}[thm]}
  \proofstepwidestar[15]{x=o}{\rA}
  \proofstepwidestar[15]{o=x}{\FormulaRefAuto{a=b\vdash b=a}[thm]{15}}
  \proofstepwidestar[1,2,8]{o\star y=o}{\FormulaRefAuto{SemigroupWithZeroLeftAbsorptionAxiom}[ax]{1,10}}
  \proofstepwidestar[1,2,8,15]{x\star y=x}{\rIE{16,17}}
  \proofstepwidestar[1,2,8,15]{x=x\star y}{\FormulaRefAuto{a=b\vdash b=a}[thm]{18}}
  \proofstepwidestar[1,2,8,15]{x\star y\in\{x,y\}}{\rIE{19,13}}
  \proofstepwidestar[21]{x=c}{\rA}
  \proofstepwidestar[22]{y=o}{\rA}
  \proofstepwidestar[22]{o=y}{\FormulaRefAuto{a=b\vdash b=a}[thm]{22}}
  \proofstepwidestar[1,2,7]{x\star o=o}{\FormulaRefAuto{SemigroupWithZeroRightAbsorptionAxiom}[ax]{1,9}}
  \proofstepwidestar[1,2,7,22]{x\star y=y}{\rIE{23,24}}
  \proofstepwidestar[1,2,7,22]{y=x\star y}{\FormulaRefAuto{a=b\vdash b=a}[thm]{25}}
  \proofstepwidestar[1,2,7,22]{x\star y\in\{x,y\}}{\rIE{26,14}}
  \proofstepwidestar[28]{y=c}{\rA}
  \proofstepwidestar[21]{c=x}{\FormulaRefAuto{a=b\vdash b=a}[thm]{21}}
  \proofstepwidestar[28]{c=y}{\FormulaRefAuto{a=b\vdash b=a}[thm]{28}}
  \proofstepwidestar[3,21]{x\star c=x}{\rIE{29,3}}
  \proofstepwidestar[3,21,28]{x\star y=x}{\rIE{30,31}}
  \proofstepwidestar[3,21,28]{x=x\star y}{\FormulaRefAuto{a=b\vdash b=a}[thm]{32}}
  \proofstepwidestar[3,21,28]{x\star y\in\{x,y\}}{\rIE{33,13}}
  \proofstepwidestar[1,2,3,7,8,21]{x\star y\in\{x,y\}}{\rOE{12,22,27,28,34}}
  \nopagebreak
  \proofstepwidestar[1,2,3,7,8]{x\star y\in\{x,y\}}{\rOE{11,15,20,21,35}}
  \nopagebreak
  \proofstepwidestar[1,2,3]{\forall x\in\{o,c\}\,\forall y\in\{o,c\}\,(x\star y\in\{x,y\})}{\rUI{\rRI{7,\rUI{\rRI{8,36}}}}}
\end{tabproofwide}

\Needspace{16\baselineskip}
\FormulaThmDeltaK[Konstantes Mengenprodukt auf dem verschwindenden Zweierquadrat]{\begin{aligned}[t]
 &\SemigroupWithZeroAlg{S}{\star}{o}\dsep c\in S\dsep c\star c=o\vdash\\[-2pt]
 &\forall X\in\PowNE{\{o,c\}}\,\forall Y\in\PowNE{\{o,c\}}\,
 (X\star_{\mathcal P}Y=\{o\})
\end{aligned}}{NPZeroPairPowerConstant}{\DeltaRow{Mengen}{S\dsep T\dsep A\dsep B\dsep P\dsep Q\dsep X\dsep Y\dsep Z\dsep K}\DeltaRow{Elemente}{o\dsep a\dsep b\dsep c\dsep d\dsep x\dsep y\dsep z\dsep u\dsep v}\DeltaRow{Funktionen}{\star\dsep\diamond\dsep F\dsep\Phi\dsep G}}
\begin{tabproofwide}
  \proofstepwidestar[1]{\SemigroupWithZeroAlg{S}{\star}{o}}{\rA}
  \proofstepwidestar[2]{c\in S}{\rA}
  \proofstepwidestar[3]{c\star c=o}{\rA}
  \proofstepwidestar[1]{\SemigroupAlg{S}{\star}}{\FormulaRefAuto{SemigroupWithZeroBaseAxiom}[ax]{1}}
  \proofstepwidestar[1]{o\in S}{\FormulaRefAuto{SemigroupWithZeroCarrierAxiom}[ax]{1}}
  \proofstepwidestar[1,2]{\{o,c\}\subseteq S}{\FormulaRefAuto{PairSetSubsetFromMembership}[thm]{5,2}}
  \proofstepwidestar[1,2,3]{\forall x\in\{o,c\}\,\forall y\in\{o,c\}\,(x\star y=o)}{\FormulaRefAuto{NPZeroPairConstant}[thm]{1,2,3}}
  \proofstepwidestar[8]{X\in\PowNE{\{o,c\}}}{\rA}
  \proofstepwidestar[9]{Y\in\PowNE{\{o,c\}}}{\rA}
  \proofstepwidestar[1,2,8]{X\in\PowNE S}{\FormulaRefAuto{NonemptyPowersetMonotone}[thm]{6,8}}
  \proofstepwidestar[1,2,9]{Y\in\PowNE S}{\FormulaRefAuto{NonemptyPowersetMonotone}[thm]{6,9}}
  \proofstepwidestar[12]{x\in X}{\rA}
  \proofstepwidestar[13]{y\in Y}{\rA}
  \proofstepwidestar[8,12]{x\in\{o,c\}}{\FormulaRefAuto{PowerSemigroupCarrierElement}[thm]{8,12}}
  \proofstepwidestar[9,13]{y\in\{o,c\}}{\FormulaRefAuto{PowerSemigroupCarrierElement}[thm]{9,13}}
  \proofstepwidestar[1,2,3,8,9,12,13]{x\star y=o}{\rRE{\rUE{\rRE{\rUE{7},14}},15}}
  \proofstepwidestar[1,2,3,8,9]{\forall x\in X\,\forall y\in Y\,(x\star y=o)}{\rUI{\rRI{12,\rUI{\rRI{13,16}}}}}
  \nopagebreak
  \proofstepwidestar[1,2,3,8,9]{X\star_{\mathcal P}Y=\{o\}}{\rRE{\rLREb{\FormulaRefAuto{PowerSemigroupConstantProduct}[thm]{4,10,11}},17}}
  \nopagebreak
  \proofstepwidestar[1,2,3]{\begin{aligned}[t]&\forall X\in\PowNE{\{o,c\}}\,\forall Y\in\PowNE{\{o,c\}}\,\\&(X\star_{\mathcal P}Y=\{o\})\end{aligned}}{\rUI{\rRI{8,\rUI{\rRI{9,18}}}}}
\end{tabproofwide}

\Needspace{16\baselineskip}
\FormulaThmDeltaK[Kommutativität der Potenzhalbgruppe des Zweierträgers]{\begin{aligned}[t]
 &\SemigroupWithZeroAlg{S}{\star}{o}\dsep c\in S\vdash\\[-2pt]
 &\forall X\in\PowNE{\{o,c\}}\,\forall Y\in\PowNE{\{o,c\}}\,
 (X\star_{\mathcal P}Y=Y\star_{\mathcal P}X)
\end{aligned}}{NPZeroPairPowerCommutative}{\DeltaRow{Mengen}{S\dsep T\dsep A\dsep B\dsep P\dsep Q\dsep X\dsep Y\dsep Z\dsep K}\DeltaRow{Elemente}{o\dsep a\dsep b\dsep c\dsep d\dsep x\dsep y\dsep z\dsep u\dsep v}\DeltaRow{Funktionen}{\star\dsep\diamond\dsep F\dsep\Phi\dsep G}}
\begin{tabproofwide}
  \proofstepwidestar[1]{\SemigroupWithZeroAlg{S}{\star}{o}}{\rA}
  \proofstepwidestar[2]{c\in S}{\rA}
  \proofstepwidestar[1]{\SemigroupAlg{S}{\star}}{\FormulaRefAuto{SemigroupWithZeroBaseAxiom}[ax]{1}}
  \proofstepwidestar[1]{o\in S}{\FormulaRefAuto{SemigroupWithZeroCarrierAxiom}[ax]{1}}
  \proofstepwidestar[1,2]{\{o,c\}\subseteq S}{\FormulaRefAuto{PairSetSubsetFromMembership}[thm]{4,2}}
  \proofstepwidestar[1,2]{\forall x\in\{o,c\}\,\forall y\in\{o,c\}\,(x\star y=y\star x)}{\FormulaRefAuto{NPZeroPairCommutative}[thm]{1,2}}
  \proofstepwidestar[7]{X\in\PowNE{\{o,c\}}}{\rA}
  \proofstepwidestar[8]{Y\in\PowNE{\{o,c\}}}{\rA}
  \proofstepwidestar[1,2,7]{X\in\PowNE S}{\FormulaRefAuto{NonemptyPowersetMonotone}[thm]{5,7}}
  \proofstepwidestar[1,2,8]{Y\in\PowNE S}{\FormulaRefAuto{NonemptyPowersetMonotone}[thm]{5,8}}
  \proofstepwidestar[11]{x\in X}{\rA}
  \proofstepwidestar[12]{y\in Y}{\rA}
  \proofstepwidestar[7,11]{x\in\{o,c\}}{\FormulaRefAuto{PowerSemigroupCarrierElement}[thm]{7,11}}
  \proofstepwidestar[8,12]{y\in\{o,c\}}{\FormulaRefAuto{PowerSemigroupCarrierElement}[thm]{8,12}}
  \proofstepwidestar[1,2,7,8,11,12]{x\star y=y\star x}{\rRE{\rUE{\rRE{\rUE{6},13}},14}}
  \proofstepwidestar[1,2,7,8]{\forall x\in X\,\forall y\in Y\,(x\star y=y\star x)}{\rUI{\rRI{11,\rUI{\rRI{12,15}}}}}
  \nopagebreak
  \proofstepwidestar[1,2,7,8]{X\star_{\mathcal P}Y=Y\star_{\mathcal P}X}{\FormulaRefAuto{NPPowerProductCommutes}[thm]{3,9,10,16}}
  \nopagebreak
  \proofstepwidestar[1,2]{\begin{aligned}[t]&\forall X\in\PowNE{\{o,c\}}\,\forall Y\in\PowNE{\{o,c\}}\,\\&(X\star_{\mathcal P}Y=Y\star_{\mathcal P}X)\end{aligned}}{\rUI{\rRI{7,\rUI{\rRI{8,17}}}}}
\end{tabproofwide}

\Needspace{16\baselineskip}
\FormulaThmDeltaK[Idempotenz der Potenzhalbgruppe im nichtverschwindenden Fall]{\begin{aligned}[t]
 &\SemigroupWithZeroAlg{S}{\star}{o}\dsep c\in S\dsep c\star c=c\vdash\\[-2pt]
 &\forall X\in\PowNE{\{o,c\}}\,(X\star_{\mathcal P}X=X)
\end{aligned}}{NPZeroPairPowerIdempotent}{\DeltaRow{Mengen}{S\dsep T\dsep A\dsep B\dsep P\dsep Q\dsep X\dsep Y\dsep Z\dsep K}\DeltaRow{Elemente}{o\dsep a\dsep b\dsep c\dsep d\dsep x\dsep y\dsep z\dsep u\dsep v}\DeltaRow{Funktionen}{\star\dsep\diamond\dsep F\dsep\Phi\dsep G}}
\begin{tabproofwide}
  \proofstepwidestar[1]{\SemigroupWithZeroAlg{S}{\star}{o}}{\rA}
  \proofstepwidestar[2]{c\in S}{\rA}
  \proofstepwidestar[3]{c\star c=c}{\rA}
  \proofstepwidestar[1]{\SemigroupAlg{S}{\star}}{\FormulaRefAuto{SemigroupWithZeroBaseAxiom}[ax]{1}}
  \proofstepwidestar[1]{o\in S}{\FormulaRefAuto{SemigroupWithZeroCarrierAxiom}[ax]{1}}
  \proofstepwidestar[1,2]{\{o,c\}\subseteq S}{\FormulaRefAuto{PairSetSubsetFromMembership}[thm]{5,2}}
  \proofstepwidestar[1,2,3]{\forall x\in\{o,c\}\,(x\star x=x)}{\FormulaRefAuto{NPZeroPairIdempotent}[thm]{1,2,3}}
  \proofstepwidestar[1,2,3]{\forall x\in\{o,c\}\,\forall y\in\{o,c\}\,(x\star y\in\{x,y\})}{\FormulaRefAuto{NPZeroPairConservative}[thm]{1,2,3}}
  \proofstepwidestar[9]{X\in\PowNE{\{o,c\}}}{\rA}
  \proofstepwidestar[1,2,9]{X\in\PowNE S}{\FormulaRefAuto{NonemptyPowersetMonotone}[thm]{6,9}}
  \proofstepwidestar[11]{x\in X}{\rA}
  \proofstepwidestar[9,11]{x\in\{o,c\}}{\FormulaRefAuto{PowerSemigroupCarrierElement}[thm]{9,11}}
  \proofstepwidestar[1,2,3,9,11]{x\star x=x}{\rRE{\rUE{7},12}}
  \proofstepwidestar[1,2,3,9]{\forall x\in X\,(x\star x=x)}{\rUI{\rRI{11,13}}}
  \proofstepwidestar[15]{y\in X}{\rA}
  \proofstepwidestar[9,15]{y\in\{o,c\}}{\FormulaRefAuto{PowerSemigroupCarrierElement}[thm]{9,15}}
  \proofstepwidestar[1,2,3,9,11,15]{x\star y\in\{x,y\}}{\rRE{\rUE{\rRE{\rUE{8},12}},16}}
  \proofstepwidestar[1,2,3,9]{\forall x\in X\,\forall y\in X\,(x\star y\in\{x,y\})}{\rUI{\rRI{11,\rUI{\rRI{15,17}}}}}
  \nopagebreak
  \proofstepwidestar[1,2,3,9]{X\star_{\mathcal P}X=X}{\FormulaRefAuto{NPPowerConservativeIdempotence}[thm]{4,10,14,18}}
  \nopagebreak
  \proofstepwidestar[1,2,3]{\forall X\in\PowNE{\{o,c\}}\,(X\star_{\mathcal P}X=X)}{\rUI{\rRI{9,19}}}
\end{tabproofwide}

\Needspace{16\baselineskip}
\FormulaThmDeltaK[Transport eines konstanten Produktgesetzes]{\begin{aligned}[t]&\SemigroupIso{F}{A,\star}{B,\diamond}\dsep d\in A\dsep \forall x\in A\,\forall y\in A\,(x\star y=d)\vdash\\&\forall u\in B\,\forall v\in B\,(u\diamond v=F(d))\end{aligned}}{NPIsoConstantLaw}{\DeltaRow{Mengen}{S\dsep T\dsep A\dsep B\dsep P\dsep Q\dsep X\dsep Y\dsep Z\dsep K}\DeltaRow{Elemente}{o\dsep a\dsep b\dsep c\dsep d\dsep x\dsep y\dsep z\dsep u\dsep v}\DeltaRow{Funktionen}{\star\dsep\diamond\dsep F\dsep\Phi\dsep G}}
\begin{tabproofwide}
  \proofstepwidestar[1]{\SemigroupIso{F}{A,\star}{B,\diamond}}{\rA}
  \proofstepwidestar[2]{d\in A}{\rA}
  \proofstepwidestar[3]{\forall x\in A\,\forall y\in A\,(x\star y=d)}{\rA}
  \proofstepwidestar[1]{F\colon A\bij B}{\FormulaRefAuto{SemigroupIsoBijectiveAxiom}[ax]{1}}
  \proofstepwidestar[1]{F^{-1}\colon B\bij A}{\FormulaRefAuto{InverseFunctionBijection}[thm]{4}}
  \proofstepwidestar[1]{F^{-1}\colon B\to A}{\FormulaRefAuto{Bijektive Funktion}[def]{5}}
  \proofstepwidestar[1]{\SemigroupHom{F}{A,\star}{B,\diamond}}{\FormulaRefAuto{SemigroupIsoHomAxiom}[ax]{1}}
  \proofstepwidestar[8]{u\in B}{\rA}
  \proofstepwidestar[1,8]{F^{-1}(u)\in A}{\FormulaRefAuto{F\colon A\to B\dsep x\in A\vdash F(x)\in B}[thm]{6,8}}
  \proofstepwidestar[1,8]{F(F^{-1}(u))=u}{\FormulaRefAuto{InverseFunctionRightInverse}[thm]{4,8}}
  \proofstepwidestar[11]{v\in B}{\rA}
  \proofstepwidestar[1,11]{F^{-1}(v)\in A}{\FormulaRefAuto{F\colon A\to B\dsep x\in A\vdash F(x)\in B}[thm]{6,11}}
  \proofstepwidestar[1,11]{F(F^{-1}(v))=v}{\FormulaRefAuto{InverseFunctionRightInverse}[thm]{4,11}}
  \proofstepwidestar[1,8,11]{\begin{aligned}[t]&F(F^{-1}(u)\star F^{-1}(v))\\&\quad=F(F^{-1}(u))\diamond F(F^{-1}(v))\end{aligned}}{\FormulaRefAuto{SemigroupHomOperationAxiom}[ax]{7,9,12}}
  \proofstepwidestar[1,8,11]{F(F^{-1}(u)\star F^{-1}(v))=u\diamond v}{\rIE{13,\rIE{10,14}}}
  \proofstepwidestar[1,3,8,11]{F^{-1}(u)\star F^{-1}(v)=d}{\rRE{\rUE{\rRE{\rUE{3},9}},12}}
  \proofstepwidestar[1,3,8,11]{F(d)=u\diamond v}{\rIE{16,15}}
  \nopagebreak
  \proofstepwidestar[1,3,8,11]{u\diamond v=F(d)}{\FormulaRefAuto{a=b\vdash b=a}[thm]{17}}
  \nopagebreak
  \proofstepwidestar[1,3]{\forall u\in B\,\forall v\in B\,(u\diamond v=F(d))}{\rUI{\rRI{8,\rUI{\rRI{11,18}}}}}
\end{tabproofwide}

\Needspace{16\baselineskip}
\FormulaThmDeltaK[Transport der Kommutativität]{\begin{aligned}[t]&\SemigroupIso{F}{A,\star}{B,\diamond}\dsep \forall x\in A\,\forall y\in A\,(x\star y=y\star x)\vdash\\&\forall u\in B\,\forall v\in B\,(u\diamond v=v\diamond u)\end{aligned}}{NPIsoCommutativeLaw}{\DeltaRow{Mengen}{S\dsep T\dsep A\dsep B\dsep P\dsep Q\dsep X\dsep Y\dsep Z\dsep K}\DeltaRow{Elemente}{o\dsep a\dsep b\dsep c\dsep d\dsep x\dsep y\dsep z\dsep u\dsep v}\DeltaRow{Funktionen}{\star\dsep\diamond\dsep F\dsep\Phi\dsep G}}
\begin{tabproofwide}
  \proofstepwidestar[1]{\SemigroupIso{F}{A,\star}{B,\diamond}}{\rA}
  \proofstepwidestar[2]{\forall x\in A\,\forall y\in A\,(x\star y=y\star x)}{\rA}
  \proofstepwidestar[1]{F\colon A\bij B}{\FormulaRefAuto{SemigroupIsoBijectiveAxiom}[ax]{1}}
  \proofstepwidestar[1]{F^{-1}\colon B\bij A}{\FormulaRefAuto{InverseFunctionBijection}[thm]{3}}
  \proofstepwidestar[1]{F^{-1}\colon B\to A}{\FormulaRefAuto{Bijektive Funktion}[def]{4}}
  \proofstepwidestar[1]{\SemigroupHom{F}{A,\star}{B,\diamond}}{\FormulaRefAuto{SemigroupIsoHomAxiom}[ax]{1}}
  \proofstepwidestar[7]{u\in B}{\rA}
  \proofstepwidestar[1,7]{F^{-1}(u)\in A}{\FormulaRefAuto{F\colon A\to B\dsep x\in A\vdash F(x)\in B}[thm]{5,7}}
  \proofstepwidestar[1,7]{F(F^{-1}(u))=u}{\FormulaRefAuto{InverseFunctionRightInverse}[thm]{3,7}}
  \proofstepwidestar[10]{v\in B}{\rA}
  \proofstepwidestar[1,10]{F^{-1}(v)\in A}{\FormulaRefAuto{F\colon A\to B\dsep x\in A\vdash F(x)\in B}[thm]{5,10}}
  \proofstepwidestar[1,10]{F(F^{-1}(v))=v}{\FormulaRefAuto{InverseFunctionRightInverse}[thm]{3,10}}
  \proofstepwidestar[1,7,10]{\begin{aligned}[t]&F(F^{-1}(u)\star F^{-1}(v))\\&\quad=F(F^{-1}(u))\diamond F(F^{-1}(v))\end{aligned}}{\FormulaRefAuto{SemigroupHomOperationAxiom}[ax]{6,8,11}}
  \proofstepwidestar[1,7,10]{F(F^{-1}(u)\star F^{-1}(v))=u\diamond v}{\rIE{12,\rIE{9,13}}}
  \proofstepwidestar[1,2,7,10]{F^{-1}(u)\star F^{-1}(v)=F^{-1}(v)\star F^{-1}(u)}{\rRE{\rUE{\rRE{\rUE{2},8}},11}}
  \proofstepwidestar[1,2,7,10]{F(F^{-1}(v)\star F^{-1}(u))=u\diamond v}{\rIE{15,14}}
  \proofstepwidestar[1,7,10]{\begin{aligned}[t]&F(F^{-1}(v)\star F^{-1}(u))\\&\quad=F(F^{-1}(v))\diamond F(F^{-1}(u))\end{aligned}}{\FormulaRefAuto{SemigroupHomOperationAxiom}[ax]{6,11,8}}
  \proofstepwidestar[1,7,10]{F(F^{-1}(v)\star F^{-1}(u))=v\diamond u}{\rIE{9,\rIE{12,17}}}
  \proofstepwidestar[1,2,7,10]{u\diamond v=F(F^{-1}(v)\star F^{-1}(u))}{\FormulaRefAuto{a=b\vdash b=a}[thm]{16}}
  \nopagebreak
  \proofstepwidestar[1,2,7,10]{u\diamond v=v\diamond u}{\rIE{18,19}}
  \nopagebreak
  \proofstepwidestar[1,2]{\forall u\in B\,\forall v\in B\,(u\diamond v=v\diamond u)}{\rUI{\rRI{7,\rUI{\rRI{10,20}}}}}
\end{tabproofwide}

\Needspace{16\baselineskip}
\FormulaThmDeltaK[Transport der Idempotenz]{\begin{aligned}[t]&\SemigroupIso{F}{A,\star}{B,\diamond}\dsep \forall x\in A\,(x\star x=x)\vdash\\&\forall u\in B\,(u\diamond u=u)\end{aligned}}{NPIsoIdempotentLaw}{\DeltaRow{Mengen}{S\dsep T\dsep A\dsep B\dsep P\dsep Q\dsep X\dsep Y\dsep Z\dsep K}\DeltaRow{Elemente}{o\dsep a\dsep b\dsep c\dsep d\dsep x\dsep y\dsep z\dsep u\dsep v}\DeltaRow{Funktionen}{\star\dsep\diamond\dsep F\dsep\Phi\dsep G}}
\begin{tabproofwide}
  \proofstepwidestar[1]{\SemigroupIso{F}{A,\star}{B,\diamond}}{\rA}
  \proofstepwidestar[2]{\forall x\in A\,(x\star x=x)}{\rA}
  \proofstepwidestar[1]{F\colon A\bij B}{\FormulaRefAuto{SemigroupIsoBijectiveAxiom}[ax]{1}}
  \proofstepwidestar[1]{F^{-1}\colon B\bij A}{\FormulaRefAuto{InverseFunctionBijection}[thm]{3}}
  \proofstepwidestar[1]{F^{-1}\colon B\to A}{\FormulaRefAuto{Bijektive Funktion}[def]{4}}
  \proofstepwidestar[1]{\SemigroupHom{F}{A,\star}{B,\diamond}}{\FormulaRefAuto{SemigroupIsoHomAxiom}[ax]{1}}
  \proofstepwidestar[7]{u\in B}{\rA}
  \proofstepwidestar[1,7]{F^{-1}(u)\in A}{\FormulaRefAuto{F\colon A\to B\dsep x\in A\vdash F(x)\in B}[thm]{5,7}}
  \proofstepwidestar[1,7]{F(F^{-1}(u))=u}{\FormulaRefAuto{InverseFunctionRightInverse}[thm]{3,7}}
  \proofstepwidestar[1,7]{\begin{aligned}[t]&F(F^{-1}(u)\star F^{-1}(u))\\&\quad=F(F^{-1}(u))\diamond F(F^{-1}(u))\end{aligned}}{\FormulaRefAuto{SemigroupHomOperationAxiom}[ax]{6,8,8}}
  \proofstepwidestar[1,7]{F(F^{-1}(u)\star F^{-1}(u))=u\diamond u}{\rIE{9,\rIE{9,10}}}
  \proofstepwidestar[1,2,7]{F^{-1}(u)\star F^{-1}(u)=F^{-1}(u)}{\rRE{\rUE{2},8}}
  \proofstepwidestar[1,2,7]{F(F^{-1}(u))=u\diamond u}{\rIE{12,11}}
  \proofstepwidestar[1,2,7]{u=u\diamond u}{\rIE{9,13}}
  \nopagebreak
  \proofstepwidestar[1,2,7]{u\diamond u=u}{\FormulaRefAuto{a=b\vdash b=a}[thm]{14}}
  \nopagebreak
  \proofstepwidestar[1,2]{\forall u\in B\,(u\diamond u=u)}{\rUI{\rRI{7,15}}}
\end{tabproofwide}

\Needspace{16\baselineskip}
\FormulaThmDeltaK[Idempotenz des Potenzträgers spiegelt sich in Einermengen]{\begin{aligned}[t]&\SemigroupAlg{T}{\diamond}\dsep T\ne\varnothing\dsep\\&\forall X\in\PowNE{T^2}\,(X\diamond_{\mathcal P}X=X)\vdash\\&\forall x\in T^2\,(x\diamond x=x)\end{aligned}}{NPPowerSquareReflectsIdempotent}{\DeltaRow{Mengen}{S\dsep T\dsep A\dsep B\dsep P\dsep Q\dsep X\dsep Y\dsep Z\dsep K}\DeltaRow{Elemente}{o\dsep a\dsep b\dsep c\dsep d\dsep x\dsep y\dsep z\dsep u\dsep v}\DeltaRow{Funktionen}{\star\dsep\diamond\dsep F\dsep\Phi\dsep G}}
\begin{tabproofwide}
  \proofstepwidestar[1]{\SemigroupAlg{T}{\diamond}}{\rA}
  \proofstepwidestar[2]{T\ne\varnothing}{\rA}
  \proofstepwidestar[3]{\forall X\in\PowNE{T^2}\,(X\diamond_{\mathcal P}X=X)}{\rA}
  \proofstepwidestar[1,2]{T^2\subseteq T}{\FormulaRefAuto{SemigroupSquareSubset}[thm]{1,2}}
  \proofstepwidestar[5]{x\in T^2}{\rA}
  \proofstepwidestar[1,2,5]{x\in T}{\FormulaRefAuto{A\subseteq B\dsep x\in A\vdash x\in B}[thm]{4,5}}
  \proofstepwidestar[5]{\{x\}\in\PowNE{T^2}}{\FormulaRefAuto{SingletonPowerMember}[thm]{5}}
  \proofstepwidestar[3,5]{\{x\}\diamond_{\mathcal P}\{x\}=\{x\}}{\rRE{\rUE{3},7}}
  \proofstepwidestar[1,2,5]{\{x\}\diamond_{\mathcal P}\{x\}=\{x\diamond x\}}{\FormulaRefAuto{SingletonPowerProduct}[thm]{1,6,6}}
  \proofstepwidestar[1,2,3,5]{\{x\diamond x\}=\{x\}}{\rIE{9,8}}
  \nopagebreak
  \proofstepwidestar[1,2,3,5]{x\diamond x=x}{\FormulaRefAuto{\{a\}=\{b\}\eqvdash a=b}[thm]{10}}
  \nopagebreak
  \proofstepwidestar[1,2,3]{\forall x\in T^2\,(x\diamond x=x)}{\rUI{\rRI{5,11}}}
\end{tabproofwide}

\Needspace{16\baselineskip}
\FormulaThmDeltaK[Kommutativität des Potenzträgers spiegelt sich in Einermengen]{\begin{aligned}[t]&\SemigroupAlg{T}{\diamond}\dsep T\ne\varnothing\dsep\\&\forall X\in\PowNE{T^2}\,\forall Y\in\PowNE{T^2}\,(X\diamond_{\mathcal P}Y=Y\diamond_{\mathcal P}X)\vdash\\&\forall x\in T^2\,\forall y\in T^2\,(x\diamond y=y\diamond x)\end{aligned}}{NPPowerSquareReflectsCommutative}{\DeltaRow{Mengen}{S\dsep T\dsep A\dsep B\dsep P\dsep Q\dsep X\dsep Y\dsep Z\dsep K}\DeltaRow{Elemente}{o\dsep a\dsep b\dsep c\dsep d\dsep x\dsep y\dsep z\dsep u\dsep v}\DeltaRow{Funktionen}{\star\dsep\diamond\dsep F\dsep\Phi\dsep G}}
\begin{tabproofwide}
  \proofstepwidestar[1]{\SemigroupAlg{T}{\diamond}}{\rA}
  \proofstepwidestar[2]{T\ne\varnothing}{\rA}
  \proofstepwidestar[3]{\begin{aligned}[t]&\forall X\in\PowNE{T^2}\,\forall Y\in\PowNE{T^2}\,\\&(X\diamond_{\mathcal P}Y=Y\diamond_{\mathcal P}X)\end{aligned}}{\rA}
  \proofstepwidestar[1,2]{T^2\subseteq T}{\FormulaRefAuto{SemigroupSquareSubset}[thm]{1,2}}
  \proofstepwidestar[5]{x\in T^2}{\rA}
  \proofstepwidestar[1,2,5]{x\in T}{\FormulaRefAuto{A\subseteq B\dsep x\in A\vdash x\in B}[thm]{4,5}}
  \proofstepwidestar[5]{\{x\}\in\PowNE{T^2}}{\FormulaRefAuto{SingletonPowerMember}[thm]{5}}
  \proofstepwidestar[8]{y\in T^2}{\rA}
  \proofstepwidestar[1,2,8]{y\in T}{\FormulaRefAuto{A\subseteq B\dsep x\in A\vdash x\in B}[thm]{4,8}}
  \proofstepwidestar[8]{\{y\}\in\PowNE{T^2}}{\FormulaRefAuto{SingletonPowerMember}[thm]{8}}
  \proofstepwidestar[3,5,8]{\{x\}\diamond_{\mathcal P}\{y\}=\{y\}\diamond_{\mathcal P}\{x\}}{\rRE{\rUE{\rRE{\rUE{3},7}},10}}
  \proofstepwidestar[1,2,5,8]{\{x\}\diamond_{\mathcal P}\{y\}=\{x\diamond y\}}{\FormulaRefAuto{SingletonPowerProduct}[thm]{1,6,9}}
  \proofstepwidestar[1,2,5,8]{\{y\}\diamond_{\mathcal P}\{x\}=\{y\diamond x\}}{\FormulaRefAuto{SingletonPowerProduct}[thm]{1,9,6}}
  \proofstepwidestar[1,2,3,5,8]{\{x\diamond y\}=\{y\diamond x\}}{\rIE{13,\rIE{12,11}}}
  \nopagebreak
  \proofstepwidestar[1,2,3,5,8]{x\diamond y=y\diamond x}{\FormulaRefAuto{\{a\}=\{b\}\eqvdash a=b}[thm]{14}}
  \nopagebreak
  \proofstepwidestar[1,2,3]{\forall x\in T^2\,\forall y\in T^2\,(x\diamond y=y\diamond x)}{\rUI{\rRI{5,\rUI{\rRI{8,15}}}}}
\end{tabproofwide}

\Needspace{16\baselineskip}
\FormulaThmDeltaK[Ein konstantes Potenzquadrat liefert eine Null des Grundquadrats]{\begin{aligned}[t]
 &\SemigroupAlg{T}{\diamond}\dsep T\ne\varnothing\dsep u\in T^2\dsep\\[-2pt]
 &\forall X\in\PowNE{T^2}\,\forall Y\in\PowNE{T^2}\,(X\diamond_{\mathcal P}Y=Z)\vdash\\[-2pt]
 &u\diamond u\in T^2\land\forall x\in T^2\,((u\diamond u)\diamond x=u\diamond u\land x\diamond(u\diamond u)=u\diamond u)
\end{aligned}}{NPPowerConstantSquareAbsorber}{\DeltaRow{Mengen}{S\dsep T\dsep A\dsep B\dsep P\dsep Q\dsep X\dsep Y\dsep Z\dsep K}\DeltaRow{Elemente}{o\dsep a\dsep b\dsep c\dsep d\dsep x\dsep y\dsep z\dsep u\dsep v}\DeltaRow{Funktionen}{\star\dsep\diamond\dsep F\dsep\Phi\dsep G}}
\begin{tabproofwide}
  \proofstepwidestar[1]{\SemigroupAlg{T}{\diamond}}{\rA}
  \proofstepwidestar[2]{T\ne\varnothing}{\rA}
  \proofstepwidestar[3]{u\in T^2}{\rA}
  \proofstepwidestar[4]{\begin{aligned}[t]&\forall X\in\PowNE{T^2}\,\forall Y\in\PowNE{T^2}\,\\&(X\diamond_{\mathcal P}Y=Z)\end{aligned}}{\rA}
  \proofstepwidestar[1,2]{T^2\subseteq T}{\FormulaRefAuto{SemigroupSquareSubset}[thm]{1,2}}
  \proofstepwidestar[1,2,3]{u\in T}{\FormulaRefAuto{A\subseteq B\dsep x\in A\vdash x\in B}[thm]{5,3}}
  \proofstepwidestar[3]{\{u\}\in\PowNE{T^2}}{\FormulaRefAuto{SingletonPowerMember}[thm]{3}}
  \proofstepwidestar[1,2,3]{u\diamond u\in T^2}{\FormulaRefAuto{SemigroupSquareProductMember}[thm]{1,2,6,6}}
  \proofstepwidestar[3,4]{\{u\}\diamond_{\mathcal P}\{u\}=Z}{\rRE{\rUE{\rRE{\rUE{4},7}},7}}
  \proofstepwidestar[1,2,3]{\{u\}\diamond_{\mathcal P}\{u\}=\{u\diamond u\}}{\FormulaRefAuto{SingletonPowerProduct}[thm]{1,6,6}}
  \proofstepwidestar[1,2,3,4]{\{u\diamond u\}=Z}{\rIE{10,9}}
  \proofstepwidestar[1,2,3,4]{Z=\{u\diamond u\}}{\FormulaRefAuto{a=b\vdash b=a}[thm]{11}}
  \proofstepwidestar[13]{x\in T^2}{\rA}
  \proofstepwidestar[14]{y\in T^2}{\rA}
  \proofstepwidestar[1,2,13]{x\in T}{\FormulaRefAuto{A\subseteq B\dsep x\in A\vdash x\in B}[thm]{5,13}}
  \proofstepwidestar[1,2,14]{y\in T}{\FormulaRefAuto{A\subseteq B\dsep x\in A\vdash x\in B}[thm]{5,14}}
  \proofstepwidestar[13]{\{x\}\in\PowNE{T^2}}{\FormulaRefAuto{SingletonPowerMember}[thm]{13}}
  \proofstepwidestar[14]{\{y\}\in\PowNE{T^2}}{\FormulaRefAuto{SingletonPowerMember}[thm]{14}}
  \proofstepwidestar[4,13,14]{\{x\}\diamond_{\mathcal P}\{y\}=Z}{\rRE{\rUE{\rRE{\rUE{4},17}},18}}
  \proofstepwidestar[1,2,13,14]{\{x\}\diamond_{\mathcal P}\{y\}=\{x\diamond y\}}{\FormulaRefAuto{SingletonPowerProduct}[thm]{1,15,16}}
  \proofstepwidestar[1,2,4,13,14]{\{x\diamond y\}=Z}{\rIE{20,19}}
  \proofstepwidestar[1,2,3,4,13,14]{\{x\diamond y\}=\{u\diamond u\}}{\rIE{12,21}}
  \proofstepwidestar[1,2,3,4,13,14]{x\diamond y=u\diamond u}{\FormulaRefAuto{\{a\}=\{b\}\eqvdash a=b}[thm]{22}}
  \proofstepwidestar[1,2,3,4]{\forall x\in T^2\,\forall y\in T^2\,(x\diamond y=u\diamond u)}{\rUI{\rRI{13,\rUI{\rRI{14,23}}}}}
  \proofstepwidestar[1,2,3,4,13]{(u\diamond u)\diamond x=u\diamond u}{\rRE{\rUE{\rRE{\rUE{24},8}},13}}
  \proofstepwidestar[1,2,3,4,13]{x\diamond(u\diamond u)=u\diamond u}{\rRE{\rUE{\rRE{\rUE{24},13}},8}}
  \proofstepwidestar[1,2,3,4,13]{(u\diamond u)\diamond x=u\diamond u\land x\diamond(u\diamond u)=u\diamond u}{\rAI{25,26}}
  \nopagebreak
  \proofstepwidestar[1,2,3,4]{\forall x\in T^2\,((u\diamond u)\diamond x=u\diamond u\land x\diamond(u\diamond u)=u\diamond u)}{\rUI{\rRI{13,27}}}
  \nopagebreak
  \proofstepwidestar[1,2,3,4]{\begin{aligned}[t]&u\diamond u\in T^2\land\\&\forall x\in T^2\,((u\diamond u)\diamond x=u\diamond u\\&\qquad{}\land x\diamond(u\diamond u)=u\diamond u)\end{aligned}}{\rAI{8,28}}
\end{tabproofwide}

\Needspace{16\baselineskip}
\FormulaThmDeltaK[Benennung der beiden Produktwerte nach Wahl der Null]{\begin{aligned}[t]
 &T^2=\{u,v\}\dsep u\ne v\dsep z\in T^2\dsep\SemigroupWithZeroAlg{T}{\diamond}{z}\vdash\\[-2pt]
 &\exists d\,(\SemigroupWithZeroAlg{T}{\diamond}{z}\land(T^2=\{z,d\}\land z\ne d))
\end{aligned}}{NPTwoElementZeroPresentation}{\DeltaRow{Mengen}{S\dsep T\dsep A\dsep B\dsep P\dsep Q\dsep X\dsep Y\dsep Z\dsep K}\DeltaRow{Elemente}{o\dsep a\dsep b\dsep c\dsep d\dsep x\dsep y\dsep z\dsep u\dsep v}\DeltaRow{Funktionen}{\star\dsep\diamond\dsep F\dsep\Phi\dsep G}}
\begin{tabproofwide}
  \proofstepwidestar[1]{T^2=\{u,v\}}{\rA}
  \proofstepwidestar[2]{u\ne v}{\rA}
  \proofstepwidestar[3]{z\in T^2}{\rA}
  \proofstepwidestar[4]{\SemigroupWithZeroAlg{T}{\diamond}{z}}{\rA}
  \proofstepwidestar[1,3]{z\in\{u,v\}}{\rIE{1,3}}
  \proofstepwidestar[1,3]{z=u\lor z=v}{\FormulaRefAuto{x\in\{A,B\}\eqvdash(x=A\lor x=B)}[thm]{5}}
  \proofstepwidestar[7]{z=u}{\rA}
  \proofstepwidestar[7]{u=z}{\FormulaRefAuto{a=b\vdash b=a}[thm]{7}}
  \proofstepwidestar[1,7]{T^2=\{z,v\}}{\rIE{8,1}}
  \proofstepwidestar[2,7]{z\ne v}{\rIE{8,2}}
  \proofstepwidestar[1,2,4,7]{\SemigroupWithZeroAlg{T}{\diamond}{z}\land(T^2=\{z,v\}\land z\ne v)}{\rAI{4,\rAI{9,10}}}
  \proofstepwidestar[1,2,4,7]{\exists d\,(\SemigroupWithZeroAlg{T}{\diamond}{z}\land(T^2=\{z,d\}\land z\ne d))}{\rEI{11}}
  \proofstepwidestar[13]{z=v}{\rA}
  \proofstepwidestar[13]{v=z}{\FormulaRefAuto{a=b\vdash b=a}[thm]{13}}
  \proofstepwidestar[]{\{u,v\}=\{v,u\}}{\FormulaRefAuto{\{a,b\}=\{b,a\}}[thm]}
  \proofstepwidestar[1]{T^2=\{v,u\}}{\rIE{15,1}}
  \proofstepwidestar[1,13]{T^2=\{z,u\}}{\rIE{14,16}}
  \proofstepwidestar[2]{v\ne u}{\FormulaRefAuto{a\neq b\vdash b\neq a}[thm]{2}}
  \proofstepwidestar[2,13]{z\ne u}{\rIE{14,18}}
  \proofstepwidestar[1,2,4,13]{\SemigroupWithZeroAlg{T}{\diamond}{z}\land(T^2=\{z,u\}\land z\ne u)}{\rAI{4,\rAI{17,19}}}
  \nopagebreak
  \proofstepwidestar[1,2,4,13]{\exists d\,(\SemigroupWithZeroAlg{T}{\diamond}{z}\land(T^2=\{z,d\}\land z\ne d))}{\rEI{20}}
  \nopagebreak
  \proofstepwidestar[1,2,3,4]{\exists d\,(\SemigroupWithZeroAlg{T}{\diamond}{z}\land(T^2=\{z,d\}\land z\ne d))}{\rOE{6,7,12,13,21}}
\end{tabproofwide}

\hypertarget{np.zero}{}
\NullproductsProofTarget{NPTargetZero}
Für die letzte Tabelle seien \(P=\PowNE S\), \(Q=\PowNE T\),
\(A=\PowNE{\{o,c\}}\), \(B=\PowNE{T^2}\) und
\(G=\Phi\restriction_{P^2}\). Dies sind Abkürzungen der angezeigten
Terme; die Operationen auf \(A,B\) sind die eingeschränkten Mengenprodukte.

\Needspace{16\baselineskip}
\FormulaThmDeltaK[Die Zielhalbgruppe besitzt eine Null und genau einen weiteren Produktwert]{\begin{aligned}[t]
 &\SemigroupWithZeroAlg{S}{\star}{o}\dsep S^2=\{o,c\}\dsep
 \SemigroupAlg{T}{\diamond}\dsep T\ne\varnothing\dsep\\[-2pt]
 &T^2=\{u,v\}\dsep u\ne v\dsep
 \SemigroupIso{\Phi}{\PowNE S,\star_{\mathcal P}}{\PowNE T,\diamond_{\mathcal P}}\vdash\\[-2pt]
 &\exists z\,\exists d\,(\SemigroupWithZeroAlg{T}{\diamond}{z}\land(T^2=\{z,d\}\land z\ne d))\end{aligned}}{NPTargetZero}{\DeltaRow{Mengen}{S\dsep T\dsep A\dsep B\dsep P\dsep Q\dsep X\dsep Y\dsep Z\dsep K}\DeltaRow{Elemente}{o\dsep a\dsep b\dsep c\dsep d\dsep x\dsep y\dsep z\dsep u\dsep v}\DeltaRow{Funktionen}{\star\dsep\diamond\dsep F\dsep\Phi\dsep G}}
\begin{tabproofwide}
  \proofstepwidestar[1]{\SemigroupWithZeroAlg{S}{\star}{o}}{\rA}
  \proofstepwidestar[2]{S^2=\{o,c\}}{\rA}
  \proofstepwidestar[3]{\SemigroupAlg{T}{\diamond}}{\rA}
  \proofstepwidestar[4]{T\ne\varnothing}{\rA}
  \proofstepwidestar[5]{T^2=\{u,v\}}{\rA}
  \proofstepwidestar[6]{u\ne v}{\rA}
  \proofstepwidestar[7]{\SemigroupIso{\Phi}{P,\star_{\mathcal P}}{Q,\diamond_{\mathcal P}}}{\rA}
  \proofstepwidestar[1]{\SemigroupAlg{S}{\star}}{\FormulaRefAuto{SemigroupWithZeroBaseAxiom}[ax]{1}}
  \proofstepwidestar[1]{o\in S}{\FormulaRefAuto{SemigroupWithZeroCarrierAxiom}[ax]{1}}
  \proofstepwidestar[1]{S\ne\varnothing}{\FormulaRefAuto{a\in A\vdash A\neq\varnothing}[thm]{9}}
  \proofstepwidestar[1]{S^2\subseteq S}{\FormulaRefAuto{SemigroupSquareSubset}[thm]{8,10}}
  \proofstepwidestar[]{c\in\{o,c\}}{\FormulaRefAuto{b\in\{a,b\}}[thm]}
  \proofstepwidestar[2]{\{o,c\}=S^2}{\FormulaRefAuto{a=b\vdash b=a}[thm]{2}}
  \proofstepwidestar[2]{c\in S^2}{\rIE{13,12}}
  \proofstepwidestar[1,2]{c\in S}{\FormulaRefAuto{A\subseteq B\dsep x\in A\vdash x\in B}[thm]{11,14}}
  \proofstepwidestar[1,2]{c\star c\in S^2}{\FormulaRefAuto{SemigroupSquareProductMember}[thm]{8,10,15,15}}
  \proofstepwidestar[1,2]{c\star c\in\{o,c\}}{\rIE{2,16}}
  \proofstepwidestar[1,2]{c\star c=o\lor c\star c=c}{\FormulaRefAuto{x\in\{A,B\}\eqvdash(x=A\lor x=B)}[thm]{17}}
  \proofstepwidestar[1]{P\ne\varnothing}{\FormulaRefAuto{NonemptyPowerCarrierNonempty}[thm]{10}}
  \proofstepwidestar[4]{Q\ne\varnothing}{\FormulaRefAuto{NonemptyPowerCarrierNonempty}[thm]{4}}
  \proofstepwidestar[1,4,7]{\SemigroupIso{G}{P^2,\star_{\mathcal P}}{Q^2,\diamond_{\mathcal P}}}{\FormulaRefAuto{SemigroupIsoSquareRestriction}[thm]{7,19,20}}
  \proofstepwidestar[1,2]{P^2=\PowNE{S^2}}{\FormulaRefAuto{PairProductSquarePowerEquality}[thm]{8,10,2}}
  \proofstepwidestar[1,2]{P^2=A}{\rIE{2,22}}
  \proofstepwidestar[3,4,5]{Q^2=B}{\FormulaRefAuto{PairProductSquarePowerEquality}[thm]{3,4,5}}
  \proofstepwidestar[1,2,3,4,5,7]{\SemigroupIso{G}{A,\star_{\mathcal P}}{B,\diamond_{\mathcal P}}}{\rIE{24,\rIE{23,21}}}
  \proofstepwidestar[]{\{o\}\in A}{\FormulaRefAuto{SingletonPowerMember}[thm]{\FormulaRefAuto{a\in\{a,b\}}[thm]}}
  \proofstepwidestar[]{u\in\{u,v\}}{\FormulaRefAuto{a\in\{a,b\}}[thm]}
  \proofstepwidestar[5]{\{u,v\}=T^2}{\FormulaRefAuto{a=b\vdash b=a}[thm]{5}}
  \proofstepwidestar[5]{u\in T^2}{\rIE{28,27}}
  \proofstepwidestar[30]{c\star c=o}{\rA}
  \proofstepwidestar[1,2,30]{\forall X\in A\,\forall Y\in A\,(X\star_{\mathcal P}Y=\{o\})}{\FormulaRefAuto{NPZeroPairPowerConstant}[thm]{1,15,30}}
  \proofstepwidestar[1,2,3,4,5,7,30]{\forall X\in B\,\forall Y\in B\,(X\diamond_{\mathcal P}Y=G(\{o\}))}{\FormulaRefAuto{NPIsoConstantLaw}[thm]{25,26,31}}
  \proofstepwidestar[1,2,3,4,5,7,30]{\begin{aligned}[t]&u\diamond u\in T^2\land\\&\forall x\in T^2\,((u\diamond u)\diamond x=u\diamond u\\&\qquad{}\land x\diamond(u\diamond u)=u\diamond u)\end{aligned}}{\FormulaRefAuto{NPPowerConstantSquareAbsorber}[thm]{3,4,29,32}}
  \proofstepwidestar[1,2,3,4,5,7,30]{\SemigroupWithZeroAlg{T}{\diamond}{u\diamond u}}{\FormulaRefAuto{NPZeroFromSquareAbsorption}[thm]{3,4,\rAEa{33},\rAEb{33}}}
  \proofstepwidestar[1,2,3,4,5,6,7,30]{\begin{aligned}[t]&\exists d\,(\SemigroupWithZeroAlg{T}{\diamond}{u\diamond u}\\&{}\land(T^2=\{u\diamond u,d\}\\&{}\land u\diamond u\ne d))\end{aligned}}{\FormulaRefAuto{NPTwoElementZeroPresentation}[thm]{5,6,\rAEa{33},34}}
  \proofstepwidestar[1,2,3,4,5,6,7,30]{\begin{aligned}[t]&\exists z\,\exists d\,(\SemigroupWithZeroAlg{T}{\diamond}{z}\\&{}\land(T^2=\{z,d\}\\&{}\land z\ne d))\end{aligned}}{\rEI{35}}
  \proofstepwidestar[37]{c\star c=c}{\rA}
  \proofstepwidestar[1,2,37]{\forall X\in A\,(X\star_{\mathcal P}X=X)}{\FormulaRefAuto{NPZeroPairPowerIdempotent}[thm]{1,15,37}}
  \proofstepwidestar[1,2]{\forall X\in A\,\forall Y\in A\,(X\star_{\mathcal P}Y=Y\star_{\mathcal P}X)}{\FormulaRefAuto{NPZeroPairPowerCommutative}[thm]{1,15}}
  \proofstepwidestar[1,2,3,4,5,7,37]{\forall X\in B\,(X\diamond_{\mathcal P}X=X)}{\FormulaRefAuto{NPIsoIdempotentLaw}[thm]{25,38}}
  \proofstepwidestar[1,2,3,4,5,7]{\forall X\in B\,\forall Y\in B\,(X\diamond_{\mathcal P}Y=Y\diamond_{\mathcal P}X)}{\FormulaRefAuto{NPIsoCommutativeLaw}[thm]{25,39}}
  \proofstepwidestar[1,2,3,4,5,7,37]{\forall x\in T^2\,(x\diamond x=x)}{\FormulaRefAuto{NPPowerSquareReflectsIdempotent}[thm]{3,4,40}}
  \proofstepwidestar[1,2,3,4,5,7]{\forall x\in T^2\,\forall y\in T^2\,(x\diamond y=y\diamond x)}{\FormulaRefAuto{NPPowerSquareReflectsCommutative}[thm]{3,4,41}}
  \proofstepwidestar[1,2,3,4,5,7,37]{\begin{aligned}[t]&u\diamond v\in T^2\land\\&\forall x\in T^2\,((u\diamond v)\diamond x=u\diamond v\\&\qquad{}\land x\diamond(u\diamond v)=u\diamond v)\end{aligned}}{\FormulaRefAuto{NPTwoElementSemilatticeAbsorber}[thm]{3,4,5,42,43}}
  \proofstepwidestar[1,2,3,4,5,7,37]{\SemigroupWithZeroAlg{T}{\diamond}{u\diamond v}}{\FormulaRefAuto{NPZeroFromSquareAbsorption}[thm]{3,4,\rAEa{44},\rAEb{44}}}
  \proofstepwidestar[1,2,3,4,5,6,7,37]{\begin{aligned}[t]&\exists d\,(\SemigroupWithZeroAlg{T}{\diamond}{u\diamond v}\\&{}\land(T^2=\{u\diamond v,d\}\\&{}\land u\diamond v\ne d))\end{aligned}}{\FormulaRefAuto{NPTwoElementZeroPresentation}[thm]{5,6,\rAEa{44},45}}
  \nopagebreak
  \proofstepwidestar[1,2,3,4,5,6,7,37]{\begin{aligned}[t]&\exists z\,\exists d\,(\SemigroupWithZeroAlg{T}{\diamond}{z}\\&{}\land(T^2=\{z,d\}\\&{}\land z\ne d))\end{aligned}}{\rEI{46}}
  \nopagebreak
  \proofstepwidestar[1,2,3,4,5,6,7]{\begin{aligned}[t]&\exists z\,\exists d\,(\SemigroupWithZeroAlg{T}{\diamond}{z}\\&{}\land(T^2=\{z,d\}\\&{}\land z\ne d))\end{aligned}}{\rOE{18,30,36,37,47}}
\end{tabproofwide}



\Needspace{26\baselineskip}
\section{Die beiden ausgezeichneten Profile}
\hypertarget{np.marked}{}
\NullproductsProofTarget{NPMarkedProfileCompatibility}

\Needspace{14\baselineskip}
\FormulaThmDeltaK[Die Null und der zweite Produktwert haben passende Profile]{%
\begin{aligned}[t]
&H_S\dsep H_T\dsep S^2=\{0_S,c\}\dsep T^2=\{0_T,d\}
 \dsep c\neq0_S\dsep I\\[-2pt]
&\vdash\Phi(\{0_S\})\sim_T\{0_T\}\land\Phi(\{c\})\sim_T\{d\}
\end{aligned}
}{NPMarkedProfileCompatibility}{%
  \DeltaRow{Mengen}{S\dsep T\dsep0_S\dsep0_T\dsep c\dsep d}
}
\begin{tabproofwide}
  \proofstepwidestar[1]{H_S}{\rA}
  \proofstepwidestar[2]{H_T}{\rA}
  \proofstepwidestar[3]{S^2=\{0_S,c\}}{\rA}
  \proofstepwidestar[4]{T^2=\{0_T,d\}}{\rA}
  \proofstepwidestar[5]{c\neq0_S}{\rA}
  \proofstepwidestar[6]{I}{\rA}
  \proofstepwidestar[1]{\SemigroupAlg{S}{\star}}{\FormulaRefAuto{SemigroupWithZeroBaseAxiom}[ax]{1}}
  \proofstepwidestar[2]{\SemigroupAlg{T}{\diamond}}{\FormulaRefAuto{SemigroupWithZeroBaseAxiom}[ax]{2}}
  \proofstepwidestar[1]{0_S\in S}{\FormulaRefAuto{SemigroupWithZeroCarrierAxiom}[ax]{1}}
  \proofstepwidestar[2]{0_T\in T}{\FormulaRefAuto{SemigroupWithZeroCarrierAxiom}[ax]{2}}
  \proofstepwidestar[1]{S\neq\varnothing}{\FormulaRefAuto{a\in A\vdash A\neq\varnothing}[thm]{9}}
  \proofstepwidestar[2]{T\neq\varnothing}{\FormulaRefAuto{a\in A\vdash A\neq\varnothing}[thm]{10}}
  \proofstepwidestar[1]{S^2\subseteq S}{\FormulaRefAuto{SemigroupSquareSubset}[thm]{7,11}}
  \proofstepwidestar[2]{T^2\subseteq T}{\FormulaRefAuto{SemigroupSquareSubset}[thm]{8,12}}
  \proofstepwidestar[3]{c\in S^2}{\rIE{\FormulaRefAuto{a=b\vdash b=a}[thm]{3},\FormulaRefAuto{b\in\{a,b\}}[thm]}}
  \proofstepwidestar[4]{d\in T^2}{\rIE{\FormulaRefAuto{a=b\vdash b=a}[thm]{4},\FormulaRefAuto{b\in\{a,b\}}[thm]}}
  \proofstepwidestar[1,3]{c\in S}{\FormulaRefAuto{A\subseteq B\dsep x\in A\vdash x\in B}[thm]{13,15}}
  \proofstepwidestar[2,4]{d\in T}{\FormulaRefAuto{A\subseteq B\dsep x\in A\vdash x\in B}[thm]{14,16}}
  \proofstepwidestar[1,3]{\{c\}\in P_S}{\FormulaRefAuto{SingletonInNonemptyPowerset}[thm]{17}}
  \proofstepwidestar[1]{\{0_S\}\in P_S}{\FormulaRefAuto{SingletonInNonemptyPowerset}[thm]{9}}
  \proofstepwidestar[2]{\{0_T\}\in P_T}{\FormulaRefAuto{SingletonInNonemptyPowerset}[thm]{10}}
  \proofstepwidestar[2,4]{\{d\}\in P_T}{\FormulaRefAuto{SingletonInNonemptyPowerset}[thm]{18}}
  \proofstepwidestar[1,2,6]{\Phi(\{0_S\})=\{0_T\}}{\FormulaRefAuto{PowerSemigroupZeroIsoValue}[thm]{1,2,6}}
  \proofstepwidestar[2]{\EqRel{P_T,\sim_T}}{\FormulaRefAuto{NullProfileEquivalence}[thm]{2}}
  \proofstepwidestar[2]{\{0_T\}\sim_T\{0_T\}}{\FormulaRefAuto{EqRelReflexivity}[ax]{24,21}}
  \proofstepwidestar[1,2,6]{\Phi(\{0_S\})\sim_T\{0_T\}}{\rIE{\FormulaRefAuto{a=b\vdash b=a}[thm]{23},25}}
  \proofstepwidestar[2,4]{\{d\}\sim_T\{d\}}{\FormulaRefAuto{EqRelReflexivity}[ax]{24,22}}
  \proofstepwidestar[2,4]{\{d\}\sim_T\{0_T,d\}}{\FormulaRefAuto{NullProfileAdjoinZero}[thm]{2,18}}
  \proofstepwidestar[2,4]{\{0_T,d\}\subseteq T}{\FormulaRefAuto{PairSetSubsetFromMembership}[thm]{10,18}}
  \proofstepwidestar[2,4]{\{0_T,d\}\in P_T}{\FormulaRefAuto{NonemptyPowersetMembership}[thm]{\rAI{29,\FormulaRefAuto{\{a,b\}\neq\varnothing}[thm]}}}
  \proofstepwidestar[2,4]{\{0_T,d\}\sim_T\{d\}}{\FormulaRefAuto{EqRelSymmetry}[ax]{24,\rAI{22,30},28}}
  \proofstepwidestar[5]{\{c\}\neq\{0_S\}}{\rAEa{\FormulaRefAuto{PairNonemptySubsetsDistinct}[thm]{5}}}
  \proofstepwidestar[1,3,6]{\{c\}=\{0_S\}\leftrightarrow\Phi(\{c\})=\Phi(\{0_S\})}{\FormulaRefAuto{SemigroupIsoEqualityReflection}[thm]{6,\rAI{19,20}}}
  \proofstepwidestar[1,3,5,6]{\Phi(\{c\})\neq\Phi(\{0_S\})}{\FormulaRefAuto{P\leftrightarrow Q,\neg P\vdash\neg Q}[thm]{33,32}}
  \proofstepwidestar[1,2,3,5,6]{\Phi(\{c\})\neq\{0_T\}}{\rIE{23,34}}
  \proofstepwidestar[1,3]{\{c\}\in P_S^2}{\rRE{\rLREb{\FormulaRefAuto{PowerSemigroupSquareSingletonCriterion}[thm]{7,11,17}},15}}
  \proofstepwidestar[1]{\SemigroupAlg{P_S}{\star_{\mathcal P}}}{\FormulaRefAuto{NonemptyPowerSemigroup}[thm]{7}}
  \proofstepwidestar[1]{P_S\neq\varnothing}{\FormulaRefAuto{NonemptyPowerCarrierNonempty}[thm]{11}}
  \proofstepwidestar[2]{P_T\neq\varnothing}{\FormulaRefAuto{NonemptyPowerCarrierNonempty}[thm]{12}}
  \proofstepwidestar[1]{P_S^2\subseteq P_S}{\FormulaRefAuto{SemigroupSquareSubset}[thm]{37,38}}
  \proofstepwidestar[6]{\Phi:P_S\to P_T}{\FormulaRefAuto{SemigroupIsoFunction}[thm]{6}}
  \proofstepwidestar[1,3,6]{\Phi(\{c\})\in\Phi[P_S^2]}{\FormulaRefAuto{C\subseteq A\dsep F\colon A\to B\dsep x\in C\vdash F(x)\in F[C]}[thm]{40,41,36}}
  \proofstepwidestar[1,2,6]{\Phi[P_S^2]=P_T^2}{\FormulaRefAuto{SemigroupIsoSquareImage}[thm]{6,38,39}}
  \proofstepwidestar[1,2,3,6]{\Phi(\{c\})\in P_T^2}{\rIE{43,42}}
  \proofstepwidestar[2]{P_T^2\subseteq\PowNE{T^2}}{\FormulaRefAuto{PowerSemigroupSquareSubset}[thm]{8,12}}
  \proofstepwidestar[1,2,3,6]{\Phi(\{c\})\in\PowNE{T^2}}{\FormulaRefAuto{A\subseteq B\dsep x\in A\vdash x\in B}[thm]{45,44}}
  \proofstepwidestar[1,2,3,6]{\Phi(\{c\})\subseteq T^2\land\Phi(\{c\})\neq\varnothing}{\FormulaRefAuto{NonemptyPowersetMembership}[thm]{46}}
  \proofstepwidestar[1,2,3,4,6]{\Phi(\{c\})\subseteq\{0_T,d\}}{\rIE{4,\rAEa{47}}}
  \proofstepwidestar[1,2,3,4,6]{\begin{aligned}[t]&\Phi(\{c\})=\{0_T\}\\&{}\lor\Phi(\{c\})=\{d\}\\&{}\lor\Phi(\{c\})=\{0_T,d\}\end{aligned}}{\FormulaRefAuto{NonemptySubsetPairCases}[thm]{48,\rAEb{47}}}
  \proofstepwidestar[1,2,3,4,5,6]{\Phi(\{c\})=\{d\}\lor\Phi(\{c\})=\{0_T,d\}}{\FormulaRefAuto{P\lor Q,\neg P\vdash Q}[thm]{49,35}}
  \proofstepwidestar[51]{\Phi(\{c\})=\{d\}}{\rA}
  \proofstepwidestar[2,4,51]{\Phi(\{c\})\sim_T\{d\}}{\rIE{\FormulaRefAuto{a=b\vdash b=a}[thm]{51},27}}
  \proofstepwidestar[53]{\Phi(\{c\})=\{0_T,d\}}{\rA}
  \proofstepwidestar[2,4,53]{\Phi(\{c\})\sim_T\{d\}}{\rIE{\FormulaRefAuto{a=b\vdash b=a}[thm]{53},31}}
  \nopagebreak
  \proofstepwidestar[1,2,3,4,5,6]{\Phi(\{c\})\sim_T\{d\}}{\rOE{50,51,52,53,54}}
  \nopagebreak
  \proofstepwidestar[1,2,3,4,5,6]{\begin{aligned}[t]&\Phi(\{0_S\})\sim_T\{0_T\}\\&{}\land\Phi(\{c\})\sim_T\{d\}\end{aligned}}{\rAI{26,55}}
\end{tabproofwide}
\Needspace{26\baselineskip}
\section{Von der Nullrelation zur Multiplikation}
\Needspace{14\baselineskip}
\FormulaThmDeltaK[Eine markierte Nullrelation genügt bei zwei Produktwerten]{%
\begin{aligned}[t]
&H_S\dsep H_T\dsep S^2=\{0_S,c\}\dsep T^2=\{0_T,d\}
 \dsep f:S\bij T\\[-2pt]
&\dsep f(0_S)=0_T\dsep f(c)=d\\[-2pt]
&\dsep\forall x\in S\,\forall y\in S\,
 (x\star y=0_S\leftrightarrow f(x)\diamond f(y)=0_T)\\[-2pt]
&\vdash\SemigroupIso{f}{S,\star}{T,\diamond}
\end{aligned}
}{NPTwoValueNullBijectionIsomorphism}{%
  \DeltaRow{Mengen}{S\dsep T\dsep0_S\dsep0_T\dsep c\dsep d\dsep x\dsep y}
}
\begin{tabproofwide}
  \proofstepwidestar[1]{H_S}{\rA}
  \proofstepwidestar[2]{H_T}{\rA}
  \proofstepwidestar[3]{S^2=\{0_S,c\}}{\rA}
  \proofstepwidestar[4]{T^2=\{0_T,d\}}{\rA}
  \proofstepwidestar[5]{f:S\bij T}{\rA}
  \proofstepwidestar[6]{f(0_S)=0_T}{\rA}
  \proofstepwidestar[7]{f(c)=d}{\rA}
  \proofstepwidestar[8]{\begin{aligned}[t]&\forall x\in S\,\forall y\in S\,\\&(x\star y=0_S\leftrightarrow f(x)\diamond f(y)=0_T)\end{aligned}}{\rA}
  \proofstepwidestar[1]{\SemigroupAlg{S}{\star}}{\FormulaRefAuto{SemigroupWithZeroBaseAxiom}[ax]{1}}
  \proofstepwidestar[2]{\SemigroupAlg{T}{\diamond}}{\FormulaRefAuto{SemigroupWithZeroBaseAxiom}[ax]{2}}
  \proofstepwidestar[1]{0_S\in S}{\FormulaRefAuto{SemigroupWithZeroCarrierAxiom}[ax]{1}}
  \proofstepwidestar[2]{0_T\in T}{\FormulaRefAuto{SemigroupWithZeroCarrierAxiom}[ax]{2}}
  \proofstepwidestar[1]{S\neq\varnothing}{\FormulaRefAuto{a\in A\vdash A\neq\varnothing}[thm]{11}}
  \proofstepwidestar[2]{T\neq\varnothing}{\FormulaRefAuto{a\in A\vdash A\neq\varnothing}[thm]{12}}
  \proofstepwidestar[5]{f:S\to T}{\FormulaRefAuto{Bijektive Funktion}[def]{5}}
  \proofstepwidestar[16]{x\in S}{\rA}
  \proofstepwidestar[17]{y\in S}{\rA}
  \proofstepwidestar[5,16]{f(x)\in T}{\FormulaRefAuto{F\colon A\to B\dsep x\in A\vdash F(x)\in B}[thm]{15,16}}
  \proofstepwidestar[5,17]{f(y)\in T}{\FormulaRefAuto{F\colon A\to B\dsep x\in A\vdash F(x)\in B}[thm]{15,17}}
  \proofstepwidestar[1,16,17]{x\star y\in S^2}{\FormulaRefAuto{SemigroupSquareProductMember}[thm]{9,13,\rAI{16,17}}}
  \proofstepwidestar[2,5,16,17]{f(x)\diamond f(y)\in T^2}{\FormulaRefAuto{SemigroupSquareProductMember}[thm]{10,14,\rAI{18,19}}}
  \proofstepwidestar[1,3,16,17]{x\star y=0_S\lor x\star y=c}{\FormulaRefAuto{x\in\{A,B\}\eqvdash(x=A\lor x=B)}[thm]{\rIE{3,20}}}
  \proofstepwidestar[2,4,5,16,17]{f(x)\diamond f(y)=0_T\lor f(x)\diamond f(y)=d}{\FormulaRefAuto{x\in\{A,B\}\eqvdash(x=A\lor x=B)}[thm]{\rIE{4,21}}}
  \proofstepwidestar[8,16,17]{x\star y=0_S\leftrightarrow f(x)\diamond f(y)=0_T}{\rRE{17,\rUE{\rRE{16,\rUE{8}}}}}
  \proofstepwidestar[]{x\star y=0_S\lor x\star y\neq0_S}{\FormulaRefAuto{P\lor\neg P}[thm]}
  \proofstepwidestar[26]{x\star y=0_S}{\rA}
  \proofstepwidestar[8,16,17,26]{f(x)\diamond f(y)=0_T}{\rRE{\rLREa{24},26}}
  \proofstepwidestar[6,26]{f(x\star y)=0_T}{\rIE{\FormulaRefAuto{a=b\vdash b=a}[thm]{26},6}}
  \proofstepwidestar[6,8,16,17,26]{f(x\star y)=f(x)\diamond f(y)}{\rIE{\FormulaRefAuto{a=b\vdash b=a}[thm]{27},28}}
  \proofstepwidestar[30]{x\star y\neq0_S}{\rA}
  \proofstepwidestar[1,3,16,17,30]{x\star y=c}{\FormulaRefAuto{P\lor Q,\neg P\vdash Q}[thm]{22,30}}
  \proofstepwidestar[8,16,17,30]{f(x)\diamond f(y)\neq0_T}{\FormulaRefAuto{P\leftrightarrow Q,\neg P\vdash\neg Q}[thm]{24,30}}
  \proofstepwidestar[2,4,5,8,16,17,30]{f(x)\diamond f(y)=d}{\FormulaRefAuto{P\lor Q,\neg P\vdash Q}[thm]{23,32}}
  \proofstepwidestar[1,3,7,16,17,30]{f(x\star y)=d}{\rIE{\FormulaRefAuto{a=b\vdash b=a}[thm]{31},7}}
  \proofstepwidestar[1,2,3,4,5,7,8,16,17,30]{f(x\star y)=f(x)\diamond f(y)}{\rIE{\FormulaRefAuto{a=b\vdash b=a}[thm]{33},34}}
  \proofstepwidestar[1,2,3,4,5,6,7,8,16,17]{f(x\star y)=f(x)\diamond f(y)}{\rOE{25,26,29,30,35}}
  \proofstepwidestar[1,2,3,4,5,6,7,8]{\forall x\in S\,\forall y\in S\,
 f(x\star y)=f(x)\diamond f(y)}{\rUI{\rRI{16,\rUI{\rRI{17,36}}}}}
  \nopagebreak
  \proofstepwidestar[1,2,3,4,5,6,7,8]{\SemigroupHom{f}{S,\star}{T,\diamond}}{\FormulaRefAuto{SemigroupHomDef}[def]{9,10,15,37}}
  \nopagebreak
  \proofstepwidestar[1,2,3,4,5,6,7,8]{\SemigroupIso{f}{S,\star}{T,\diamond}}{\FormulaRefAuto{SemigroupIsoDef}[def]{38,5}}
\end{tabproofwide}

\Needspace{24\baselineskip}
\section{Zusammenführung zum Rekonstruktionssatz}
\NullproductsProofTarget{NPReconstructionConclusion}
Für die folgende Tabelle verwenden wir die Abkürzung \(\mathcal K(f)\)
aus \FormulaRefAuto{NullProfileTwoPointReconstruction}[thm]:
Sie enthält \(f:S\bij T\) und die Verträglichkeit der Singletonprofile.
Für den letzten Schritt wird nur der erste Bestandteil benötigt.

\Needspace{14\baselineskip}
\FormulaThmDeltaK[Der Isomorphismus der Grundhalbgruppen]{%
\begin{aligned}[t]
&\Finite S\dsep H_S\dsep S^2=\{0_S,c\}\dsep c\neq0_S
 \dsep\SemigroupAlg{T}{\diamond}\dsep I\\[-2pt]
&\vdash\exists f\,\SemigroupIso{f}{S,\star}{T,\diamond}
\end{aligned}
}{NPReconstructionConclusion}{%
  \DeltaRow{Mengen}{S\dsep T\dsep0_S\dsep c}
}
\begin{tabproofwide}
  \proofstepwidestar[1]{\Finite S}{\rA}
  \proofstepwidestar[2]{H_S}{\rA}
  \proofstepwidestar[3]{S^2=\{0_S,c\}}{\rA}
  \proofstepwidestar[4]{c\neq0_S}{\rA}
  \proofstepwidestar[5]{\SemigroupAlg{T}{\diamond}}{\rA}
  \proofstepwidestar[6]{I}{\rA}
  \proofstepwidestar[1,2,6]{\Finite T\land T\neq\varnothing}{\FormulaRefAuto{NPFiniteNonemptyTarget}[thm]{1,2,6}}
  \proofstepwidestar[1,2,3,4,5,6]{\begin{aligned}[t]&\exists u\,\exists v\,(u\neq v\land T^2=\{u,v\}\\&{}\land P_T^2=\PowNE{T^2})\end{aligned}}{\FormulaRefAuto{NPTwoProductValuesTransport}[thm]{2,3,4,5,\rAEb{7},6}}
  \proofstepwidestar[9]{\begin{aligned}[t]&\exists v\,(u\neq v\land T^2=\{u,v\}\\&{}\land P_T^2=\PowNE{T^2})\end{aligned}}{\rA}
  \proofstepwidestar[10]{\begin{aligned}[t]&u\neq v\land T^2=\{u,v\}\\&{}\land P_T^2=\PowNE{T^2}\end{aligned}}{\rA}
  \proofstepwidestar[1,2,3,5,6,10]{\begin{aligned}[t]&\exists z\,\exists d\,(\SemigroupWithZeroAlg{T}{\diamond}{z}\\&{}\land(T^2=\{z,d\}\\&{}\land z\neq d))\end{aligned}}{\FormulaRefAuto{NPTargetZero}[thm]{2,3,5,\rAEb{7},\rAEn{10},\rAEa{10},6}}
  \proofstepwidestar[12]{\exists d\,(H_T\land(T^2=\{0_T,d\}\land0_T\neq d))}{\rA}
  \proofstepwidestar[13]{H_T\land(T^2=\{0_T,d\}\land0_T\neq d)}{\rA}
  \proofstepwidestar[13]{H_T}{\rAEa{13}}
  \proofstepwidestar[13]{T^2=\{0_T,d\}}{\rAEn{13}}
  \proofstepwidestar[13]{0_T\neq d}{\rAEn{13}}
  \proofstepwidestar[2,3,4,6,13]{\begin{aligned}[t]&\Phi(\{0_S\})\sim_T\{0_T\}\\&{}\land\Phi(\{c\})\sim_T\{d\}\end{aligned}}{\FormulaRefAuto{NPMarkedProfileCompatibility}[thm]{2,14,3,15,4,6}}
  \proofstepwidestar[2]{\SemigroupAlg{S}{\star}}{\FormulaRefAuto{SemigroupWithZeroBaseAxiom}[ax]{2}}
  \proofstepwidestar[2]{0_S\in S}{\FormulaRefAuto{SemigroupWithZeroCarrierAxiom}[ax]{2}}
  \proofstepwidestar[13]{0_T\in T}{\FormulaRefAuto{SemigroupWithZeroCarrierAxiom}[ax]{14}}
  \proofstepwidestar[2]{S\neq\varnothing}{\FormulaRefAuto{a\in A\vdash A\neq\varnothing}[thm]{19}}
  \proofstepwidestar[2]{S^2\subseteq S}{\FormulaRefAuto{SemigroupSquareSubset}[thm]{18,21}}
  \proofstepwidestar[1,2,5,6]{T^2\subseteq T}{\FormulaRefAuto{SemigroupSquareSubset}[thm]{5,\rAEb{7}}}
  \proofstepwidestar[3]{c\in S^2}{\rIE{\FormulaRefAuto{a=b\vdash b=a}[thm]{3},\FormulaRefAuto{b\in\{a,b\}}[thm]}}
  \proofstepwidestar[13]{d\in T^2}{\rIE{\FormulaRefAuto{a=b\vdash b=a}[thm]{15},\FormulaRefAuto{b\in\{a,b\}}[thm]}}
  \proofstepwidestar[2,3]{c\in S}{\FormulaRefAuto{A\subseteq B\dsep x\in A\vdash x\in B}[thm]{22,24}}
  \proofstepwidestar[1,2,5,6,13]{d\in T}{\FormulaRefAuto{A\subseteq B\dsep x\in A\vdash x\in B}[thm]{23,25}}
  \proofstepwidestar[4]{0_S\neq c}{\FormulaRefAuto{a\neq b\vdash b\neq a}[thm]{4}}
  \proofstepwidestar[1,2,3,4,5,6,13]{\begin{aligned}[t]&\exists f\,\bigl(\mathcal K(f)\land f(0_S)=0_T\land f(c)=d\\&{}\land \forall x,y\in S\,\\&\quad(x\star y=0_S\leftrightarrow f(x)\diamond f(y)=0_T)\bigr)\end{aligned}}{\FormulaRefAuto{NullProfileTwoPointReconstruction}[thm]{2,14,6,1,\rAEa{7},\rAI{19,26},\rAI{20,27},28,16,\rAEa{17},\rAEb{17}}}
  \proofstepwidestar[30]{\begin{aligned}[t]&\mathcal K(f)\land f(0_S)=0_T\land f(c)=d\\&{}\land \forall x,y\in S\,\\&\quad(x\star y=0_S\leftrightarrow f(x)\diamond f(y)=0_T)\end{aligned}}{\rA}
  \proofstepwidestar[30]{f:S\bij T}{\rAEa{\rAEa{30}}}
  \proofstepwidestar[2,3,13,30]{\SemigroupIso{f}{S,\star}{T,\diamond}}{\FormulaRefAuto{NPTwoValueNullBijectionIsomorphism}[thm]{2,14,3,15,31,\rAEn{30},\rAEn{30},\rAEn{30}}}
  \proofstepwidestar[2,3,13,30]{\exists f\,\SemigroupIso{f}{S,\star}{T,\diamond}}{\rEI{32}}
  \proofstepwidestar[1,2,3,4,5,6,13]{\exists f\,\SemigroupIso{f}{S,\star}{T,\diamond}}{\rEE{29,30,33}}
  \proofstepwidestar[1,2,3,4,5,6,12]{\exists f\,\SemigroupIso{f}{S,\star}{T,\diamond}}{\rEE{12,13,34}}
  \proofstepwidestar[1,2,3,4,5,6,10]{\exists f\,\SemigroupIso{f}{S,\star}{T,\diamond}}{\rEE{11,12,35}}
  \nopagebreak
  \proofstepwidestar[1,2,3,4,5,6,9]{\exists f\,\SemigroupIso{f}{S,\star}{T,\diamond}}{\rEE{9,10,36}}
  \nopagebreak
  \proofstepwidestar[1,2,3,4,5,6]{\exists f\,\SemigroupIso{f}{S,\star}{T,\diamond}}{\rEE{8,9,37}}
\end{tabproofwide}

Damit ist \FormulaRefAuto{FiniteTwoProductZeroReconstruction}[thm] bewiesen.
Die Endlichkeit wurde ausschließlich zur Rekonstruktion der Häufigkeiten
der Profile verwendet. Die algebraische Bestimmung der Zielnull benötigt
keine Endlichkeitsannahme.

